
\Data\ID{1003108904}
\Updated{2021-08-17 11:08:22}
\Level{A}
\SubArea{Applications of definite integral}
\LANGen{
\Question{
What is the area of the plain region  bounded by the lines $x=2$, $x=3$ and the graphs of functions $f(x)=\sin x$ and $g(x)=\frac1{x^2}$? Round the result to three decimal digits.}
{$0.407$}{$0.167$}{$1.573$}{$0.623$}{}{}}
\LANGpl{
\Question{
Jakie jest  pole powierzchni obszaru ograniczonego prostymi  $x=2$, $x=3$ oraz wykresami funkcji  $f(x)=\sin x$ i  $g(x)=\frac1{x^2}$? Zaokrągli wynik do trzech miejsc po przecinku.}
{$0{,}407$}{$0{,}167$}{$1{,}573$}{$0{,}623$}{}{}}
\LANGcs{
\Question{
Jaký obsah má plocha ohraničená přímkami $x=2$, $x=3$ a grafy funkcí $f(x)=\sin x$ a $g(x)=\frac1{x^2}$? Zaokrouhli na 3 desetinná místa.}
{$0{,}407$}{$0{,}167$}{$1{,}573$}{$0{,}623$}{}{}}
\LANGsk{
\Question{
Aký obsah má plocha ohraničená priamkami $x=2$, $x=3$ a grafy funkcí $f(x)=\sin x$ a $g(x)=\frac1{x^2}$? Zaokrúhli na 3 desatinné miesta.}
{$0{,}407$}{$0{,}167$}{$1{,}573$}{$0{,}623$}{}{}}
\LANGes{
\Question{
¿Cuál es el área de la figura plana limitada por las rectas $x=2$, $x=3$ y las gráficas de funciones $f(x)=\sin x$ y $g(x)=\frac1{x^2}$? Aproxima el resultado a tres cifras decimales.}
{$0.407$}{$0.167$}{$1.573$}{$0.623$}{}{}}
\End

\Data\ID{1003108905}
\Updated{2020-07-15 03:07:12}
\Level{A}
\SubArea{Applications of definite integral}
\LANGen{
\Question{
Find the area of the plain region bounded by the curves $f(x)=x^5$ and $g(x)=x^9$ on the interval $[1;2]$.}
{$91.8$}{$113.2$}{$91.6$}{$100.4$}{}{}}
\LANGpl{
\Question{
Oblicz pole powierzchni obszaru ograniczonego krzywymi  $f(x)=x^5$ i $g(x)=x^9$ w przedziale $\langle1;2\rangle$.}
{$91{,}8$}{$113{,}2$}{$91{,}6$}{$100{,}4$}{}{}}
\LANGcs{
\Question{
Určete obsah plochy vymezené křivkami $f(x)=x^5$ a $g(x)=x^9$ na intervalu $\langle1;2\rangle$.}
{$91{,}8$}{$113{,}2$}{$91{,}6$}{$100{,}4$}{}{}}
\LANGsk{
\Question{
Určte obsah plochy vymedzenej krivkami $f(x)=x^5$ a $g(x)=x^9$ na intervale $\langle1;2\rangle$.}
{$91{,}8$}{$113{,}2$}{$91{,}6$}{$100{,}4$}{}{}}
\LANGes{
\Question{
Encuentra el área de la superficie limitada por las curvas $f(x)=x^5$ y $g(x)=x^9$ en el intervalo $[1;2]$.}
{$91.8$}{$113.2$}{$91.6$}{$100.4$}{}{}}
\End

\Data\ID{1003108907}
\Updated{2020-07-15 03:07:12}
\Level{A}
\SubArea{Applications of definite integral}
\LANGen{
\Question{
What are the lengths of a right triangle legs, if the triangle area is given by the expression $\int\limits_{-4}^1(0.8x+4.2)\,\mathrm{d}x-5$?}
{$5$ and $4$}{$6$ and $3$}{$4$ and $6$}{$6$ and $5$}{}{}}
\LANGpl{
\Question{
Jaka jest długość ramion trójkąta prostokątnego, jeśli pole powierzchni trójkąta wyrażone jest  $\int\limits_{-4}^1(0{,}8x+4{,}2)\,\mathrm{d}x-5$?}
{$5$ i $4$}{$6$ i $3$}{$4$ i $6$}{$6$ i $5$}{}{}}
\LANGcs{
\Question{
Jak dlouhé odvěsny má pravoúhlý trojúhelník, jehož obsah je dán výrazem  $\int\limits_{-4}^1(0{,}8x+4{,}2)\,\mathrm{d}x-5$?}
{$5$ a $4$}{$6$ a $3$}{$4$ a $6$}{$6$ a $5$}{}{}}
\LANGsk{
\Question{
Aké dlhé odvesny má pravouhlý trojuholník, ktorého obsah je daný výrazom  $\int\limits_{-4}^1(0{,}8x+4{,}2)\,\mathrm{d}x-5$?}
{$5$ a $4$}{$6$ a $3$}{$4$ a $6$}{$6$ a $5$}{}{}}
\LANGes{
\Question{
¿Cuánto miden los catetos de un triángulo rectángulo si el área del triángulo está definida por $\int\limits_{-4}^1(0.8x+4.2)\,\mathrm{d}x-5$?}
{$5$ y $4$}{$6$ y $3$}{$4$ y $6$}{$6$ y $5$}{}{}}
\End

\Data\ID{1003108908}
\Updated{2021-08-17 11:08:54}
\Level{A}
\SubArea{Applications of definite integral}
\LANGen{
\Question{
What is the radius of a circle inscribed in a square, if the area of the square is given by the expression $\int\limits_1^6(6-x)\,\mathrm{d}x$ ?}
{$\frac{5\sqrt2}4$}{$3\sqrt2$}{$2\sqrt2$}{$\frac{7\sqrt2}2$}{}{}}
\LANGpl{
\Question{
Jaki jest promień okręgu wpisanego w kwadrat, pole powierzchni kwardatu jest wyrażone równaniem $\int\limits_1^6(6-x)\,\mathrm{d}x$ ?}
{$\frac{5\sqrt2}4$}{$3\sqrt2$}{$2\sqrt2$}{$\frac{7\sqrt2}2$}{}{}}
\LANGcs{
\Question{
Jaký poloměr má kružnice vepsaná čtverci, jehož obsah je dán výrazem  $\int\limits_1^6(6-x)\,\mathrm{d}x$ ?}
{$\frac{5\sqrt2}4$}{$3\sqrt2$}{$2\sqrt2$}{$\frac{7\sqrt2}2$}{}{}}
\LANGsk{
\Question{
Aký polomer má kružnica vpísaná do štvorca, ktorého obsah je daný výrazom  $\int\limits_1^6(6-x)\,\mathrm{d}x$ ?}
{$\frac{5\sqrt2}4$}{$3\sqrt2$}{$2\sqrt2$}{$\frac{7\sqrt2}2$}{}{}}
\LANGes{
\Question{
¿Cuál es el radio del círculo inscrito en un cuadrado si el área del cuadrado está definida por la expresión $\int\limits_1^6(6-x)\,\mathrm{d}x$ ?}
{$\frac{5\sqrt2}4$}{$3\sqrt2$}{$2\sqrt2$}{$\frac{7\sqrt2}2$}{}{}}
\End

\Data\ID{1003124406}
\Updated{2021-08-17 11:08:19}
\Level{A}
\SubArea{Applications of definite integral}
\LANGen{
\Question{
Calculate the area of a plane region bounded by the graphs of the functions $f(x)=x^2+2x+2$ and $g(x)=6-x^2$.}
{$9$}{$21$}{$15$}{$\frac{59}3$}{}{}}
\LANGpl{
\Question{
Oblicz  pole powierzchni obszaru ograniczonego wykresami funkcji  $f(x)=x^2+2x+2$ i $g(x)=6-x^2$.}
{$9$}{$21$}{$15$}{$\frac{59}3$}{}{}}
\LANGcs{
\Question{
Jaký obsah má rovinný obrazec, který je vymezen grafy funkcí  $f(x)=x^2+2x+2$ a $g(x)=6-x^2$?}
{$9$}{$21$}{$15$}{$\frac{59}3$}{}{}}
\LANGsk{
\Question{
Aký obsah má rovinný obrazec, ktorý je vymedzený grafmi funkcií  $f(x)=x^2+2x+2$ a $g(x)=6-x^2$?}
{$9$}{$21$}{$15$}{$\frac{59}3$}{}{}}
\LANGes{
\Question{
Calcula el área de una figura plana limitada por las gráficas de las funciones $f(x)=x^2+2x+2$ y $g(x)=6-x^2$.}
{$9$}{$21$}{$15$}{$\frac{59}3$}{}{}}
\End

\Data\ID{1003124410}
\Updated{2021-08-17 11:08:22}
\Level{A}
\SubArea{Applications of definite integral}
\LANGen{
\Question{
Calculate the area of a plane region bounded by the graphs of the functions $f(x)=-\frac x{\pi}+1$ and $g(x)=-\sin x$, and by the lines $x=\pi$ and $x=0$.}
{$2+\frac{\pi}2$}{$2-\frac{\pi}2$}{$-2-\frac{\pi}2$}{$\frac{2+\pi}2$}{}{}}
\LANGpl{
\Question{
Oblicz pole powierzchni obszaru ograniczonego wykresami funkcji  $f(x)=-\frac x{\pi}+1$ i $g(x)=-\sin x$ oraz prostymi $x=\pi$ i $x=0$.}
{$2+\frac{\pi}2$}{$2-\frac{\pi}2$}{$-2-\frac{\pi}2$}{$\frac{2+\pi}2$}{}{}}
\LANGcs{
\Question{
Jaký obsah má rovinný obrazec vymezený grafy funkcí $f(x)=-\frac x{\pi}+1$, $g(x)=-\sin x$ a přímkami $x=\pi$ and $x=0$?}
{$2+\frac{\pi}2$}{$2-\frac{\pi}2$}{$-2-\frac{\pi}2$}{$\frac{2+\pi}2$}{}{}}
\LANGsk{
\Question{
Aký obsah má rovinný obrazec vymedzený grafmi funkcií $f(x)=-\frac x{\pi}+1$, $g(x)=-\sin x$ a priamkami $x=\pi$ and $x=0$?}
{$2+\frac{\pi}2$}{$2-\frac{\pi}2$}{$-2-\frac{\pi}2$}{$\frac{2+\pi}2$}{}{}}
\LANGes{
\Question{
Calcula el área de una figura plana limitada por las gráficas de las funciones $f(x)=-\frac x{\pi}+1$ y $g(x)=-\sin x$, y por las rectas $x=\pi$ y $x=0$.}
{$2+\frac{\pi}2$}{$2-\frac{\pi}2$}{$-2-\frac{\pi}2$}{$\frac{2+\pi}2$}{}{}}
\End

\Data\ID{1103067901}
\Updated{2020-07-15 03:07:12}
\Level{A}
\SubArea{Applications of definite integral}
\IMG{1103067901_0.pdf}
\LANGen{
\Question{
Find the area of the yellow region indicated in the picture.}
{\( 6 \)}{\( 0 \)}{\( 3 \)}{\( -6 \)}{}{}}
\LANGpl{
\Question{
Oblicz powierzchnię obszaru zaznaczonego na rysunku żółtym kolorem.}
{\( 6 \)}{\( 0 \)}{\( 3 \)}{\( -6 \)}{}{}}
\LANGcs{
\Question{
Vypočítejte obsah plochy žlutě vyznačené na obrázku.}
{\( 6 \)}{\( 0 \)}{\( 3 \)}{\( -6 \)}{}{}}
\LANGsk{
\Question{
Vypočítajte obsah plochy žlto vyznačenej na obrázku.}
{\( 6 \)}{\( 0 \)}{\( 3 \)}{\( -6 \)}{}{}}
\LANGes{
\Question{
Calcula el área de la superficie amarilla marcada en la imagen.}
{\( 6 \)}{\( 0 \)}{\( 3 \)}{\( -6 \)}{}{}}
\End

\Data\ID{1103067902}
\Updated{2020-07-15 03:07:12}
\Level{A}
\SubArea{Applications of definite integral}
\IMG{1103067902.pdf}
\LANGen{
\Question{
Find the area of the region  bounded by the curves \( y=x^3-5x^2+3x+5\text{, }x=0\text{, }x=1\text{, }y=-x+5 \) as indicated in the picture.}
{\( \frac7{12} \)}{\( \frac{41}{12} \)}{\( \frac{47}{12} \)}{\( \frac5{12} \)}{}{}}
\LANGpl{
\Question{
Oblicz powierzchnię obszaru ograniczonego  krzywymi \( y=x^3-5x^2+3x+5\text{, }x=0\text{, }x=1\text{, }y=-x+5 \) wskazanymi na wykresie.}
{\( \frac7{12} \)}{\( \frac{41}{12} \)}{\( \frac{47}{12} \)}{\( \frac5{12} \)}{}{}}
\LANGcs{
\Question{
Vypočítejte obsah plochy ohraničené grafy funkcí \( y=x^3-5x^2+3x+5\text{, }x=0\text{, }x=1\text{, }y=-x+5 \) vykreslených na obrázku.}
{\( \frac7{12} \)}{\( \frac{41}{12} \)}{\( \frac{47}{12} \)}{\( \frac5{12} \)}{}{}}
\LANGsk{
\Question{
Vypočítajte obsah plochy ohraničenej grafmi funkcií \( y=x^3-5x^2+3x+5\text{, }x=0\text{, }x=1\text{, }y=-x+5 \) vykreslených na obrázku.}
{\( \frac7{12} \)}{\( \frac{41}{12} \)}{\( \frac{47}{12} \)}{\( \frac5{12} \)}{}{}}
\LANGes{
\Question{
Encuentra el área de la superficie limitada por las curvas \( y=x^3-5x^2+3x+5\text{, }x=0\text{, }x=1\text{, }y=-x+5 \) como se representa en la imagen.}
{\( \frac7{12} \)}{\( \frac{41}{12} \)}{\( \frac{47}{12} \)}{\( \frac5{12} \)}{}{}}
\End

\Data\ID{1103067903}
\Updated{2020-07-15 03:07:12}
\Level{A}
\SubArea{Applications of definite integral}
\IMG{1103067903_0.pdf}
\LANGen{
\Question{
Find the area of the yellow  region indicated in the picture.}
{\( 8\pi \)}{\( 6\pi \)}{\( 2\pi \)}{\( 10\pi \)}{}{}}
\LANGpl{
\Question{
Oblicz powierzchnię obszaru zaznaczonego na rysunku żółtym kolorem.}
{\( 8\pi \)}{\( 6\pi \)}{\( 2\pi \)}{\( 10\pi \)}{}{}}
\LANGcs{
\Question{
Vypočítej obsah plochy žlutě vyznačené na obrázku.}
{\( 8\pi \)}{\( 6\pi \)}{\( 2\pi \)}{\( 10\pi \)}{}{}}
\LANGsk{
\Question{
Vypočítaj obsah plochy žlto vyznačenej na obrázku.}
{\( 8\pi \)}{\( 6\pi \)}{\( 2\pi \)}{\( 10\pi \)}{}{}}
\LANGes{
\Question{
Calcula el área de la superficie amarilla marcada en la imagen.}
{\( 8\pi \)}{\( 6\pi \)}{\( 2\pi \)}{\( 10\pi \)}{}{}}
\End

\Data\ID{1103068001}
\Updated{2020-07-15 03:07:21}
\Level{A}
\SubArea{Applications of definite integral}
\IMG{1103068001.pdf}
\LANGen{
\Question{
Which of the following formulas does NOT express the area of the yellow triangle indicated in the picture?}
{\( \int\limits_1^ 6(-0.8x+5.8)\,\mathrm{d}x \)}{\( \frac12\cdot(5-1)\cdot(6-1)\cdot\sin90^{\circ} \)}{\( \frac{4\cdot5}2 \)}{\( \int\limits_1^ 6(-0.8x+5.8)\,\mathrm{d}x-5 \)}{}{}}
\LANGpl{
\Question{
Która z poniższych formuł NIE określa obszaru żółtego trójkąta pokazanego na obrazku?}
{\( \int\limits_1^ 6(-0{,}8x+5{,}8)\,\mathrm{d}x \)}{\( \frac12\cdot(5-1)\cdot(6-1)\cdot\sin90^{\circ} \)}{\( \frac{4\cdot5}2 \)}{\( \int\limits_1^ 6(-0{,}8x+5{,}8)\,\mathrm{d}x-5 \)}{}{}}
\LANGcs{
\Question{
Který z uvedených výrazů nevyjadřuje obsah žlutě vyznačeného trojúhelníku?}
{\( \int\limits_1^ 6(-0{,}8x+5{,}8)\,\mathrm{d}x \)}{\( \frac12\cdot(5-1)\cdot(6-1)\cdot\sin90^{\circ} \)}{\( \frac{4\cdot5}2 \)}{\( \int\limits_1^ 6(-0{,}8x+5{,}8)\,\mathrm{d}x-5 \)}{}{}}
\LANGsk{
\Question{
Ktorý z uvedených výrazov nevyjadruje obsah žlto vyznačeného trojuholníka?}
{\( \int\limits_1^ 6(-0{,}8x+5{,}8)\,\mathrm{d}x \)}{\( \frac12\cdot(5-1)\cdot(6-1)\cdot\sin90^{\circ} \)}{\( \frac{4\cdot5}2 \)}{\( \int\limits_1^ 6(-0{,}8x+5{,}8)\,\mathrm{d}x-5 \)}{}{}}
\LANGes{
\Question{
¿Cuál de las siguiente fórmulas NO expresa el área del triángulo amarillo representado en la imagen?}
{\( \int\limits_1^ 6(-0.8x+5.8)\,\mathrm{d}x \)}{\( \frac12\cdot(5-1)\cdot(6-1)\cdot\sin90^{\circ} \)}{\( \frac{4\cdot5}2 \)}{\( \int\limits_1^ 6(-0.8x+5.8)\,\mathrm{d}x-5 \)}{}{}}
\End

\Data\ID{1103068002}
\Updated{2021-08-17 11:08:46}
\Level{A}
\SubArea{Applications of definite integral}
\IMG{1103068002_0.pdf}
\LANGen{
\Question{
Define an unknown real positive constant \( a \) so that the yellow surface indicated in the picture has an area of \( 9 \) square units.}
{\( a=3 \)}{\( a=27 \)}{\( a=9 \)}{\( a=1 \)}{}{}}
\LANGpl{
\Question{
Wskaż stałą  \( a \) tak, aby pole powierzchni  żółtego  obszaru pokazanego  na rysunku  wynosiło \( 9 \) jednostek kwadratowych}
{\( a=3 \)}{\( a=27 \)}{\( a=9 \)}{\( a=1 \)}{}{}}
\LANGcs{
\Question{
Určete kladnou reálnou neznámou \( a \) tak, aby měla žlutě vybarvená plocha obsah  \( 9 \) čtverečných jednotek.}
{\( a=3 \)}{\( a=27 \)}{\( a=9 \)}{\( a=1 \)}{}{}}
\LANGsk{
\Question{
Určte kladnú reálnu neznámu \( a \) tak, aby mala žlto vyfarbená plocha obsah  \( 9 \) štvorcových jednotiek.}
{\( a=3 \)}{\( a=27 \)}{\( a=9 \)}{\( a=1 \)}{}{}}
\LANGes{
\Question{
Calcula el valor de la constante real positiva \( a \) para que el área de la superficie amarilla marcada en la imagen sea de \( 9 \) unidades cuadradas.}
{\( a=3 \)}{\( a=27 \)}{\( a=9 \)}{\( a=1 \)}{}{}}
\End

\Data\ID{1103068003}
\Updated{2021-08-19 10:08:22}
\Level{A}
\SubArea{Applications of definite integral}
\IMG{1103068003.pdf}
\LANGen{
\Question{
Find the missing real positive constant \( a \) so that the area of the yellow triangle indicated in the picture is \( 12 \) square units.}
{\( a=\frac23 \)}{\( a=\frac43 \)}{\( a=1 \)}{The constant \( a \) can not be determined from the picture.}{}{}}
\LANGpl{
\Question{
Wskaż brakującą  stałą  \( a \) tak, aby pole powierzchni żółtego trójkąta znajdującego się na rysunku było równe \( 12 \) jednostek kwadratowych.}
{\( a=\frac23 \)}{\( a=\frac43 \)}{\( a=1 \)}{Stała  \( a \) nie może być określona na podstawie rysunku.}{}{}}
\LANGcs{
\Question{
Určete chybějící kladnou reálnou konstantu  \( a \) tak, aby měl vyznačený žlutý trojúhelník obsah \( 12 \) čtverečných jednotek.}
{\( a=\frac23 \)}{\( a=\frac43 \)}{\( a=1 \)}{Konstantu \( a \) nelze z obrázku určit.}{}{}}
\LANGsk{
\Question{
Určte chýbajúcu kladnú reálnu konštantu  \( a \) tak, aby mal vyznačený žltý trojuholník obsah \( 12 \) štvorcových jednotiek.}
{\( a=\frac23 \)}{\( a=\frac43 \)}{\( a=1 \)}{Konštanta \( a \) sa nedá z obrázku určiť.}{}{}}
\LANGes{
\Question{
Encuentra la constante real positiva \( a \) de manera que el área del triángulo amarillo mostrado en la imagen sea \( 12 \) unidades cuadradas.}
{\( a=\frac23 \)}{\( a=\frac43 \)}{\( a=1 \)}{The constant \( a \) can not be determined from the picture.}{}{}}
\End

\Data\ID{1103068004}
\Updated{2021-08-19 10:08:39}
\Level{A}
\SubArea{Applications of definite integral}
\IMG{1103068004.pdf}
\LANGen{
\Question{
Find the missing real constant \( a \) so that the ratio of the green and the red area indicated in the picture is \( 4:1 \).}
{\( a=-\frac{5}3\pi \)}{\( a=-2\pi \)}{\( a=-\pi \)}{\( a=-\frac{5}4\pi \)}{}{}}
\LANGpl{
\Question{
Znajdź brakującą stałą rzeczywistą \( a \) tak,  aby stosunek zielonego do czerwonego obszaru wskazanego na obrazie wynosił  \( 4:1 \).}
{\( a=-\frac{5}3\pi \)}{\( a=-2\pi \)}{\( a=-\pi \)}{\( a=-\frac{5}4\pi \)}{}{}}
\LANGcs{
\Question{
Najděte chybějící reálnou konstantu \( a \) tak, aby byl poměr obsahů zeleně a červeně vyznačené plochy  \( 4:1 \).}
{\( a=-\frac{5}3\pi \)}{\( a=-2\pi \)}{\( a=-\pi \)}{\( a=-\frac{5}4\pi \)}{}{}}
\LANGsk{
\Question{
Nájdite chýbajúcu reálnu konštantu \( a \) tak, aby bol pomer obsahov zelenej a červenej vyznačenej plochy  \( 4:1 \).}
{\( a=-\frac{5}3\pi \)}{\( a=-2\pi \)}{\( a=-\pi \)}{\( a=-\frac{5}4\pi \)}{}{}}
\LANGes{
\Question{
Encuentra la constante real  \( a \) para que la proporción entre el área verde y  la roja marcadas en la imagen sea de \( 4:1 \).}
{\( a=-\frac{5}3\pi \)}{\( a=-2\pi \)}{\( a=-\pi \)}{\( a=-\frac{5}4\pi \)}{}{}}
\End

\Data\ID{1103068005}
\Updated{2021-08-19 10:08:36}
\Level{A}
\SubArea{Applications of definite integral}
\IMG{1103068005.pdf}
\LANGen{
\Question{
Find the missing real constant \( a \) so that the green area and the red area indicated in the picture do equal.}
{\( a=-2\pi \)}{\( a=-\frac32\pi \)}{\( a=-\frac{\pi}2 \)}{\( a=-3\pi \)}{}{}}
\LANGpl{
\Question{
Wskaż stałą rzeczywistą  \( a \) tak,  aby zielony obszar i czerwony obszar pokazany na rysunku były równe}
{\( a=-2\pi \)}{\( a=-\frac32\pi \)}{\( a=-\frac{\pi}2 \)}{\( a=-3\pi \)}{}{}}
\LANGcs{
\Question{
Určete chybějící reálnou konstantu \( a \) tak, aby byly obsahy červeně a zeleně vyznačené plochy stejné.}
{\( a=-2\pi \)}{\( a=-\frac32\pi \)}{\( a=-\frac{\pi}2 \)}{\( a=-3\pi \)}{}{}}
\LANGsk{
\Question{
Určte chýbajúcu reálnu konštantu \( a \) tak, aby boli obsahy červeno a zeleno vyznačenej plochy rovnaké.}
{\( a=-2\pi \)}{\( a=-\frac32\pi \)}{\( a=-\frac{\pi}2 \)}{\( a=-3\pi \)}{}{}}
\LANGes{
\Question{
Encuentra la constanta real \( a \) para que las áreas marcadas en rojo y verde sean iguales.}
{\( a=-2\pi \)}{\( a=-\frac32\pi \)}{\( a=-\frac{\pi}2 \)}{\( a=-3\pi \)}{}{}}
\End

\Data\ID{1103068101}
\Updated{2021-08-15 09:08:51}
\Level{A}
\SubArea{Applications of definite integral}
\IMG{1103068101.pdf}
\LANGen{
\Question{
Find the area of the yellow region that is bounded by the parabola and the line as is indicated in the picture.   Read all the needed values from the picture.}
{\( 4.5 \)}{\( 22.5 \)}{\( 3.75 \)}{\( 10.25 \)}{}{}}
\LANGpl{
\Question{
Oblicz powierzchnię obszaru zaznaczonego żółtym kolorem ograniczonym przez parabolę i prostą. Odczytaj potrzebne dane z wykresu.}
{\( 4{,}5 \)}{\( 22{,}5 \)}{\( 3{,}75 \)}{\( 10{,}25 \)}{}{}}
\LANGcs{
\Question{
Určete obsah žlutě vyznačené plochy vymezené na obrázku parabolou a přímkou. Potřebné hodnoty vyčtěte z grafu.}
{\( 4{,}5 \)}{\( 22{,}5 \)}{\( 3{,}75 \)}{\( 10{,}25 \)}{}{}}
\LANGsk{
\Question{
Určte obsah žlto vyznačenej plochy vymedzenej na obrázku parabolou a priamkou. Potrebné hodnoty vyčítate z grafu.}
{\( 4{,}5 \)}{\( 22{,}5 \)}{\( 3{,}75 \)}{\( 10{,}25 \)}{}{}}
\LANGes{
\Question{
Calcula el área de la superficie amarilla limitada por la parábola y la recta (mira la imagen). Lee todos los valores necesarios de la imagen.}
{\( 4.5 \)}{\( 22.5 \)}{\( 3.75 \)}{\( 10.25 \)}{}{}}
\End

\Data\ID{1103068102}
\Updated{2020-07-15 03:07:12}
\Level{A}
\SubArea{Applications of definite integral}
\IMG{1103068102.pdf}
\LANGen{
\Question{
Find the area of the yellow triangle ABC indicated in the picture.  Read all the needed values from the picture.}
{\( 10 \)}{\( 11 \)}{\( 9.5 \)}{\( 10.5 \)}{}{}}
\LANGpl{
\Question{
Oblicz powierzchnię trójkąta  ABC.  Odczytaj potrzebne dane z wykresu.}
{\( 10 \)}{\( 11 \)}{\( 9{,}5 \)}{\( 10{,}5 \)}{}{}}
\LANGcs{
\Question{
Určete obsah žlutě vyznačeného trojúhelníku ACB. Potřebné hodnoty vyčtěte z grafu.}
{\( 10 \)}{\( 11 \)}{\( 9{,}5 \)}{\( 10{,}5 \)}{}{}}
\LANGsk{
\Question{
Určte obsah žlto vyznačeného trojuholníka ACB. Potrebné hodnoty vyčítate z grafu.}
{\( 10 \)}{\( 11 \)}{\( 9{,}5 \)}{\( 10{,}5 \)}{}{}}
\LANGes{
\Question{
Calcula el área del triángulo amarillo ABC marcado en la imagen. Lee todos los valores necesarios de la imagen.}
{\( 10 \)}{\( 11 \)}{\( 9.5 \)}{\( 10.5 \)}{}{}}
\End

\Data\ID{1103068103}
\Updated{2020-07-15 03:07:12}
\Level{A}
\SubArea{Applications of definite integral}
\IMG{1103068103_0.pdf}
\LANGen{
\Question{
Find the area of the yellow  region indicated in the picture.}
{\( 2\sqrt2 \)}{\( 2\sqrt3 \)}{\( \sqrt2 \)}{\( 4\sqrt2 \)}{}{}}
\LANGpl{
\Question{
Oblicz powierzchnię obszaru zaznaczonego na rysunku żółtym kolorem.}
{\( 2\sqrt2 \)}{\( 2\sqrt3 \)}{\( \sqrt2 \)}{\( 4\sqrt2 \)}{}{}}
\LANGcs{
\Question{
Určete obsah žlutě vyznačené plochy.}
{\( 2\sqrt2 \)}{\( 2\sqrt3 \)}{\( \sqrt2 \)}{\( 4\sqrt2 \)}{}{}}
\LANGsk{
\Question{
Určte obsah žlto vyznačenej plochy.}
{\( 2\sqrt2 \)}{\( 2\sqrt3 \)}{\( \sqrt2 \)}{\( 4\sqrt2 \)}{}{}}
\LANGes{
\Question{
Encuentra el área de la superficie amarilla representada en la imagen.}
{\( 2\sqrt2 \)}{\( 2\sqrt3 \)}{\( \sqrt2 \)}{\( 4\sqrt2 \)}{}{}}
\End

\Data\ID{1103108901}
\Updated{2020-07-15 03:07:12}
\Level{A}
\SubArea{Applications of definite integral}
\IMG{1103108901.pdf}
\LANGen{
\Question{
Which of the given expressions does not describe the area of the region highlighted in yellow? (See the picture.)}
{$4\cdot\int\limits_0^{\frac{\pi}4}4\sin x\,\mathrm{d}x$}{$4\cdot\int\limits_0^{\pi}\sin x\,\mathrm{d}x$}{$8\cdot\int\limits_0^{\pi}\frac{\sin x}2\,\mathrm{d}x$}{$\int\limits_0^{\frac{\pi}2}16\cdot\frac{\sin x}2\,\mathrm{d}x$}{}{}}
\LANGpl{
\Question{
Który z podanych wyrażeń nie opisuje pola powierzchni obszaru zaznaczonego kolorem żółtym? (Spójrz na rysunek.)}
{$4\cdot\int\limits_0^{\frac{\pi}4}4\sin x\,\mathrm{d}x$}{$4\cdot\int\limits_0^{\pi}\sin x\,\mathrm{d}x$}{$8\cdot\int\limits_0^{\pi}\frac{\sin x}2\,\mathrm{d}x$}{$\int\limits_0^{\frac{\pi}2}16\cdot\frac{\sin x}2\,\mathrm{d}x$}{}{}}
\LANGcs{
\Question{
Kterým z uvedených výrazů nelze vyjádřit obsah žlutě vyznačené plochy?}
{$4\cdot\int\limits_0^{\frac{\pi}4}4\sin x\,\mathrm{d}x$}{$4\cdot\int\limits_0^{\pi}\sin x\,\mathrm{d}x$}{$8\cdot\int\limits_0^{\pi}\frac{\sin x}2\,\mathrm{d}x$}{$\int\limits_0^{\frac{\pi}2}16\cdot\frac{\sin x}2\,\mathrm{d}x$}{}{}}
\LANGsk{
\Question{
Ktorým z uvedených výrazov sa nedá vyjadriť obsah žlto vyznačenej plochy?}
{$4\cdot\int\limits_0^{\frac{\pi}4}4\sin x\,\mathrm{d}x$}{$4\cdot\int\limits_0^{\pi}\sin x\,\mathrm{d}x$}{$8\cdot\int\limits_0^{\pi}\frac{\sin x}2\,\mathrm{d}x$}{$\int\limits_0^{\frac{\pi}2}16\cdot\frac{\sin x}2\,\mathrm{d}x$}{}{}}
\LANGes{
\Question{
¿Cuál de las expresiones dadas no describe el área de la superficie amarilla? (Mira la imagen.)}
{$4\cdot\int\limits_0^{\frac{\pi}4}4\sin x\,\mathrm{d}x$}{$4\cdot\int\limits_0^{\pi}\sin x\,\mathrm{d}x$}{$8\cdot\int\limits_0^{\pi}\frac{\sin x}2\,\mathrm{d}x$}{$\int\limits_0^{\frac{\pi}2}16\cdot\frac{\sin x}2\,\mathrm{d}x$}{}{}}
\End

\Data\ID{1103108902}
\Updated{2020-07-15 03:07:12}
\Level{A}
\SubArea{Applications of definite integral}
\IMG{1103108902.pdf}
\LANGen{
\Question{
Which of the given expressions does not describe the area of the region highlighted in yellow? (See the picture.)}
{$\int\limits_1^3\frac1x\,\mathrm{d}x$}{$2\int\limits_1^3\frac1x\,\mathrm{d}x$}{$\int\limits_1^4\frac2x\,\mathrm{d}x-\int\limits_3^4\frac2x\,\mathrm{d}x$}{$\int\limits_1^3\frac1x\,\mathrm{d}x-\int\limits_1^3-\frac1x\,\mathrm{d}x$}{}{}}
\LANGpl{
\Question{
Które z podanych wyrażeń nie opisuje pola powierzchni obszaru zaznaczonego kolorem żółtym. (Spójrz na rysunek.)}
{$\int\limits_1^3\frac1x\,\mathrm{d}x$}{$2\int\limits_1^3\frac1x\,\mathrm{d}x$}{$\int\limits_1^4\frac2x\,\mathrm{d}x-\int\limits_3^4\frac2x\,\mathrm{d}x$}{$\int\limits_1^3\frac1x\,\mathrm{d}x-\int\limits_1^3-\frac1x\,\mathrm{d}x$}{}{}}
\LANGcs{
\Question{
Kterým z uvedených výrazů nelze vyjádřit obsah žlutě vyznačené plochy?}
{$\int\limits_1^3\frac1x\,\mathrm{d}x$}{$2\int\limits_1^3\frac1x\,\mathrm{d}x$}{$\int\limits_1^4\frac2x\,\mathrm{d}x-\int\limits_3^4\frac2x\,\mathrm{d}x$}{$\int\limits_1^3\frac1x\,\mathrm{d}x-\int\limits_1^3-\frac1x\,\mathrm{d}x$}{}{}}
\LANGsk{
\Question{
Ktorým z uvedených výrazov sa nedá vyjadriť obsah žlto vyznačenej plochy?}
{$\int\limits_1^3\frac1x\,\mathrm{d}x$}{$2\int\limits_1^3\frac1x\,\mathrm{d}x$}{$\int\limits_1^4\frac2x\,\mathrm{d}x-\int\limits_3^4\frac2x\,\mathrm{d}x$}{$\int\limits_1^3\frac1x\,\mathrm{d}x-\int\limits_1^3-\frac1x\,\mathrm{d}x$}{}{}}
\LANGes{
\Question{
¿Cuál de las expresiones no describe el área de la región sombreada en amarillo? (Mira la imagen.)}
{$\int\limits_1^3\frac1x\,\mathrm{d}x$}{$2\int\limits_1^3\frac1x\,\mathrm{d}x$}{$\int\limits_1^4\frac2x\,\mathrm{d}x-\int\limits_3^4\frac2x\,\mathrm{d}x$}{$\int\limits_1^3\frac1x\,\mathrm{d}x-\int\limits_1^3-\frac1x\,\mathrm{d}x$}{}{}}
\End

\Data\ID{1103108903}
\Updated{2020-07-15 03:07:10}
\Level{A}
\SubArea{Applications of definite integral}
\LANGen{
\Question{
Which of the given pictures shows the shaded region of the biggest area?}
{}{}{}{}{}{}}
\LANGpl{
\Question{
Na którym z rysunków znajduje się zacieniony obszar z największą powierzchnią?}
{}{}{}{}{}{}}
\LANGcs{
\Question{
Na kterém obrázku má žlutě vyznačená plocha největší obsah?}
{}{}{}{}{}{}}
\LANGsk{
\Question{
Na ktorom obrázku má žlto vyznačená plocha najväčší obsah?}
{}{}{}{}{}{}}
\LANGes{
\Question{
¿Cuál de las imágenes representa el área sombreada más grande?}
{}{}{}{}{}{}}

\IMGanswer{1103108903a_0.pdf}
\IMGanswer{1103108903b_0.pdf}
\IMGanswer{1103108903c_0.pdf}
\IMGanswer{1103108903d_0.pdf}\End

\Data\ID{1103108906}
\Updated{2021-08-15 09:08:39}
\Level{A}
\SubArea{Applications of definite integral}
\IMG{1103108906.pdf}
\LANGen{
\Question{
Find the area of the triangle $ACB$ highlighted in yellow (see the picture).}
{$20.5$}{$22.5$}{$20$}{$21$}{}{}}
\LANGpl{
\Question{
Wskaż pole powierzchni trójkąta $ACB$ oznaczonego żółtym kolorem (spójrz na rysunek).}
{$20{,}5$}{$22{,}5$}{$20$}{$21$}{}{}}
\LANGcs{
\Question{
Určete obsah žlutě vyznačeného trojúhelníku $ACB$.}
{$20{,}5$}{$22{,}5$}{$20$}{$21$}{}{}}
\LANGsk{
\Question{
Určte obsah žlto vyznačeného trojuholníka $ACB$.}
{$20{,}5$}{$22{,}5$}{$20$}{$21$}{}{}}
\LANGes{
\Question{
Encuentra el área del triángulo $ACB$ marcado en color amarillo (mira la imagen).}
{$20.5$}{$22.5$}{$20$}{$21$}{}{}}
\End

\Data\ID{1103124401}
\Updated{2021-08-15 09:08:50}
\Level{A}
\SubArea{Applications of definite integral}
\IMG{1103124401_0.pdf}
\LANGen{
\Question{
Find the area of the region highlighted in red (see the picture).}
{$1.5\pi$}{$2\pi$}{$2.5\pi$}{$4.5\pi$}{}{}}
\LANGpl{
\Question{
Oblicz pole powierzchni obszaru znaczonego czerwonym kolorem. ( Spójrz na rysunek)}
{$1{,}5\pi$}{$2\pi$}{$2{,}5\pi$}{$4{,}5\pi$}{}{}}
\LANGcs{
\Question{
Určete obsah červeně vyznačené plochy.}
{$1{,}5\pi$}{$2\pi$}{$2{,}5\pi$}{$4{,}5\pi$}{}{}}
\LANGsk{
\Question{
Určte obsah červeno vyznačenej plochy.}
{$1{,}5\pi$}{$2\pi$}{$2{,}5\pi$}{$4{,}5\pi$}{}{}}
\LANGes{
\Question{
Encuentra el área de la superficie marcada en color rojo (mira la imagen).}
{$1.5\pi$}{$2\pi$}{$2.5\pi$}{$4.5\pi$}{}{}}
\End

\Data\ID{9000065601}
\Updated{2021-08-19 10:08:15}
\Level{A}
\SubArea{Applications of definite integral}
\LANGen{
\Question{
Find the area of the region bounded by
\(x\)-axis,
graph of \(f(x) = x + 3\) 
and lines \(x = -1\) 
and \(x = 1\).}
{\(6\)}{\(2\)}{\(4\)}{\(8\)}{}{}}
\LANGpl{
\Question{
Oblicz pole obszaru ograniczonego osią 
\(x\),
wykresem  funkcji \(f\colon y = x + 3\) oraz
prostymi  \(x = -1\) 
i \(x = 1\).}
{\(6\)}{\(2\)}{\(4\)}{\(8\)}{}{}}
\LANGcs{
\Question{
Vypočítejte obsah plochy ohraničené osou
\(x\), grafem
funkce \(f\colon y = x + 3\) a
přímkami \(x = -1\) 
a \(x = 1\).}
{\(6\)}{\(2\)}{\(4\)}{\(8\)}{}{}}
\LANGsk{
\Question{
Vypočítajte obsah plochy ohraničenej osou
\(x\), grafom
funkcie \(f\colon y = x + 3\) a
priamkami \(x = -1\) 
a \(x = 1\).}
{\(6\)}{\(2\)}{\(4\)}{\(8\)}{}{}}
\LANGes{
\Question{
Encuentra el área limitada por el eje \(x\),
la gráfica de \(f(x) = x + 3\) 
y las rectas \(x = -1\) y \(x = 1\).}
{\(6\)}{\(2\)}{\(4\)}{\(8\)}{}{}}
\End

\Data\ID{9000065602}
\Updated{2021-08-19 10:08:58}
\Level{A}
\SubArea{Applications of definite integral}
\LANGen{
\Question{
Find the area of the region bounded by
\(x\)-axis,
graph of \(f(x)= x^{2} + 3\) 
and lines \(x = -2\) 
and \(x = 1\).}
{\(12\)}{\(6\)}{\(8\)}{\(10\)}{}{}}
\LANGpl{
\Question{
Oblicz pole obszaru ograniczonego osią
\(x\),
wykresem funkcji  \(f\colon y = x^{2} + 3\) 
oraz prostymi  \(x = -2\) 
i \(x = 1\).}
{\(12\)}{\(6\)}{\(8\)}{\(10\)}{}{}}
\LANGcs{
\Question{
Vypočítejte obsah plochy ohraničené osou
\(x\), grafem
funkce \(f\colon y = x^{2} + 3\) a
křivkami \(x = -2\) 
a \(x = 1\).}
{\(12\)}{\(6\)}{\(8\)}{\(10\)}{}{}}
\LANGsk{
\Question{
Vypočítajte obsah plochy ohraničenej osou
\(x\), grafom
funkcie \(f\colon y = x^{2} + 3\) a
krivkami \(x = -2\) 
a \(x = 1\).}
{\(12\)}{\(6\)}{\(8\)}{\(10\)}{}{}}
\LANGes{
\Question{
Encuentra el área de la región limitada por el eje \(x\), la gráfica de \(f(x)= x^{2} + 3\) 
y las rectas \(x = -2\) 
y \(x = 1\).}
{\(12\)}{\(6\)}{\(8\)}{\(10\)}{}{}}
\End

\Data\ID{9000065603}
\Updated{2020-07-15 03:07:37}
\Level{A}
\SubArea{Applications of definite integral}
\LANGen{
\Question{
Find the area of the region bounded by the curves
\(y = 0\),
\(y = x^{3}\),
\(x = 1\) and
\(x = 3\).}
{\(20\)}{\(22\)}{\(24\)}{\(26\)}{}{}}
\LANGpl{
\Question{
Oblicz  pole obszaru ograniczonego krzywymi
\(y = 0\),
\(y = x^{3}\),
\(x = 1\) i
\(x = 3\).}
{\(20\)}{\(22\)}{\(24\)}{\(26\)}{}{}}
\LANGcs{
\Question{
Vypočítejte obsah plochy ohraničené křivkami:
\(y = 0\),
\(y = x^{3}\),
\(x = 1\),
\(x = 3\).}
{\(20\)}{\(22\)}{\(24\)}{\(26\)}{}{}}
\LANGsk{
\Question{
Vypočítajte obsah plochy ohraničenej krivkami:
\(y = 0\),
\(y = x^{3}\),
\(x = 1\),
\(x = 3\).}
{\(20\)}{\(22\)}{\(24\)}{\(26\)}{}{}}
\LANGes{
\Question{
Encuentra el área de la superficie limitada por las curvas
\(y = 0\),
\(y = x^{3}\),
\(x = 1\) y
\(x = 3\).}
{\(20\)}{\(22\)}{\(24\)}{\(26\)}{}{}}
\End

\Data\ID{9000065604}
\Updated{2020-07-22 07:07:02}
\Level{A}
\SubArea{Applications of definite integral}
\LANGen{
\Question{
Find the area of the region bounded by the graph of
\(f(x)=\cos x\) on
\(\left [ \frac{\pi }{2};\pi \right ] \) and
lines \(y = 0\) 
and \(x =\pi \).}
{\(1\)}{\(\frac{3}
{4}\)}{\(\frac{\sqrt{3}}
 {2} \)}{\(2\)}{}{}}
\LANGpl{
\Question{
Oblicz  pole obszaru ograniczonego wykresem funkcji
\(f\colon y =\cos x\) na 
\(\left \langle \frac{\pi }{2};\pi \right \rangle \) i prostymi
 \(y = 0\) 
i \(x =\pi \).}
{\(1\)}{\(\frac{3}
{4}\)}{\(\frac{\sqrt{3}}
 {2} \)}{\(2\)}{}{}}
\LANGcs{
\Question{
Vypočítejte obsah plochy ohraničené grafem funkce
\(f\colon y =\cos x\),
\(D(f) = \left \langle \frac{\pi }{2};\pi \right \rangle \) a
přímkami \(y = 0\) 
a \(x =\pi \).}
{\(1\)}{\(\frac{3}
{4}\)}{\(\frac{\sqrt{3}}
 {2} \)}{\(2\)}{}{}}
\LANGsk{
\Question{
Vypočítajte obsah plochy ohraničenej grafom funkcie
\(f\colon y =\cos x\),
\(D(f) = \left \langle \frac{\pi }{2};\pi \right \rangle \) a
priamkami \(y = 0\) 
a \(x =\pi \).}
{\(1\)}{\(\frac{3}
{4}\)}{\(\frac{\sqrt{3}}
 {2} \)}{\(2\)}{}{}}
\LANGes{
\Question{
Encuentra el área de la región limitada por la gráfica de \(f(x)=\cos x\) en
\(\left [ \frac{\pi }{2};\pi \right ] \) y
las rectas \(y = 0\) 
y \(x =\pi \).}
{\(1\)}{\(\frac{3}
{4}\)}{\(\frac{\sqrt{3}}
 {2} \)}{\(2\)}{}{}}
\End

\Data\ID{9000065605}
\Updated{2020-07-15 03:07:37}
\Level{A}
\SubArea{Applications of definite integral}
\LANGen{
\Question{
Find the area of the region bounded by the curves
\(y = -2x\) and
\(y = -x^{2} + 3\).}
{\(\frac{32}
 {3} \)}{\(\frac{29}
 {3} \)}{\(\frac{31}
 {3} \)}{\(\frac{35}
 {3} \)}{}{}}
\LANGpl{
\Question{
Oblicz  pole obszaru ograniczonego krzywymi
\(y = -2x\) i
\(y = -x^{2} + 3\).}
{\(\frac{32}
 {3} \)}{\(\frac{29}
 {3} \)}{\(\frac{31}
 {3} \)}{\(\frac{35}
 {3} \)}{}{}}
\LANGcs{
\Question{
Vypočítejte obsah plochy ohraničené křivkami:
\(y = -2x\),
\(y = -x^{2} + 3\).}
{\(\frac{32}
 {3} \)}{\(\frac{29}
 {3} \)}{\(\frac{31}
 {3} \)}{\(\frac{35}
 {3} \)}{}{}}
\LANGsk{
\Question{
Vypočítajte obsah plochy ohraničenej krivkami:
\(y = -2x\),
\(y = -x^{2} + 3\).}
{\(\frac{32}
 {3} \)}{\(\frac{29}
 {3} \)}{\(\frac{31}
 {3} \)}{\(\frac{35}
 {3} \)}{}{}}
\LANGes{
\Question{
Encuentra el área de la superficie limitada por las curvas
\(y = -2x\) y
\(y = -x^{2} + 3\).}
{\(\frac{32}
 {3} \)}{\(\frac{29}
 {3} \)}{\(\frac{31}
 {3} \)}{\(\frac{35}
 {3} \)}{}{}}
\End

\Data\ID{9000065606}
\Updated{2020-07-15 03:07:37}
\Level{A}
\SubArea{Applications of definite integral}
\LANGen{
\Question{
Find the area of the region bounded by the curves
\(y =\mathrm{e} ^{x}\),
\(y = -\mathrm{e}^{x} + 2\) and
\(x = -3\).}
{\(4 + \frac{2} 
{\mathrm{e}^{3}} \)}{\(4 + \frac{1} 
{\mathrm{e}^{3}} \)}{\(4 -\frac{2} 
{\mathrm{e}^{3}} \)}{\(4 -\frac{1} 
{\mathrm{e}^{3}} \)}{}{}}
\LANGpl{
\Question{
Oblicz  pole obszaru ograniczonego krzywymi
\(y =\mathrm{e} ^{x}\),
\(y = -\mathrm{e}^{x} + 2\) i
\(x = -3\).}
{\(4 + \frac{2} 
{\mathrm{e}^{3}} \)}{\(4 + \frac{1} 
{\mathrm{e}^{3}} \)}{\(4 -\frac{2} 
{\mathrm{e}^{3}} \)}{\(4 -\frac{1} 
{\mathrm{e}^{3}} \)}{}{}}
\LANGcs{
\Question{
Vypočítejte obsah plochy ohraničené křivkami:
\(y =\mathrm{e} ^{x}\),
\(y = -\mathrm{e}^{x} + 2\),
\(x = -3\).}
{\(4 + \frac{2} 
{\mathrm{e}^{3}} \)}{\(4 + \frac{1} 
{\mathrm{e}^{3}} \)}{\(4 -\frac{2} 
{\mathrm{e}^{3}} \)}{\(4 -\frac{1} 
{\mathrm{e}^{3}} \)}{}{}}
\LANGsk{
\Question{
Vypočítajte obsah plochy ohraničenej krivkami:
\(y =\mathrm{e} ^{x}\),
\(y = -\mathrm{e}^{x} + 2\),
\(x = -3\).}
{\(4 + \frac{2} 
{\mathrm{e}^{3}} \)}{\(4 + \frac{1} 
{\mathrm{e}^{3}} \)}{\(4 -\frac{2} 
{\mathrm{e}^{3}} \)}{\(4 -\frac{1} 
{\mathrm{e}^{3}} \)}{}{}}
\LANGes{
\Question{
Encuentra el área de la superficie limitada por las curvas 
\(y =\mathrm{e} ^{x}\),
\(y = -\mathrm{e}^{x} + 2\) y
\(x = -3\).}
{\(4 + \frac{2} 
{\mathrm{e}^{3}} \)}{\(4 + \frac{1} 
{\mathrm{e}^{3}} \)}{\(4 -\frac{2} 
{\mathrm{e}^{3}} \)}{\(4 -\frac{1} 
{\mathrm{e}^{3}} \)}{}{}}
\End

\Data\ID{9000065607}
\Updated{2020-07-15 03:07:37}
\Level{A}
\SubArea{Applications of definite integral}
\LANGen{
\Question{
Find the area of the region bounded by the curves
\(y = 3^{x}\),
\(y = 3^{-x}\) and
\(y = 3\).}
{\(6 -\frac{4} 
{\ln 3}\)}{\(3 -\frac{2} 
{\ln 3}\)}{\(3 + \frac{4} 
{\ln 3}\)}{\(6 -\frac{2} 
{\ln 3}\)}{}{}}
\LANGpl{
\Question{
Oblicz  pole obszaru ograniczonego krzywymi
\(y = 3^{x}\),
\(y = 3^{-x}\) and
\(y = 3\).}
{\(6 -\frac{4} 
{\ln 3}\)}{\(3 -\frac{2} 
{\ln 3}\)}{\(3 + \frac{4} 
{\ln 3}\)}{\(6 -\frac{2} 
{\ln 3}\)}{}{}}
\LANGcs{
\Question{
Vypočítejte obsah plochy ohraničené křivkami:
\(y = 3^{x}\),
\(y = 3^{-x}\),
\(y = 3\).}
{\(6 -\frac{4} 
{\ln 3}\)}{\(3 -\frac{2} 
{\ln 3}\)}{\(3 + \frac{4} 
{\ln 3}\)}{\(6 -\frac{2} 
{\ln 3}\)}{}{}}
\LANGsk{
\Question{
Vypočítajte obsah plochy ohraničenej krivkami:
\(y = 3^{x}\),
\(y = 3^{-x}\),
\(y = 3\).}
{\(6 -\frac{4} 
{\ln 3}\)}{\(3 -\frac{2} 
{\ln 3}\)}{\(3 + \frac{4} 
{\ln 3}\)}{\(6 -\frac{2} 
{\ln 3}\)}{}{}}
\LANGes{
\Question{
Encuentra el área de la superficie limitada por las curvas
\(y = 3^{x}\),
\(y = 3^{-x}\),
\(y = 3\).}
{\(6 -\frac{4} 
{\ln 3}\)}{\(3 -\frac{2} 
{\ln 3}\)}{\(3 + \frac{4} 
{\ln 3}\)}{\(6 -\frac{2} 
{\ln 3}\)}{}{}}
\End

\Data\ID{9000065608}
\Updated{2020-07-22 07:07:54}
\Level{A}
\SubArea{Applications of definite integral}
\IMG{000656_08-figure0.pdf}
\LANGen{
\Question{
Using integrals write formula for the area of the shaded region.}
{\(\int _{a}^{b}(f(x) - g(x))\, \mathrm{d}x +\int _{ b}^{c}(g(x) - f(x))\, \mathrm{d}x\)}{\(\int _{a}^{b}(g(x) - f(x))\, \mathrm{d}x +\int _{ b}^{c}(g(x) - f(x))\, \mathrm{d}x\)}{\(\int _{a}^{b}(f(x) - g(x))\, \mathrm{d}x +\int _{ b}^{c}(f(x) - g(x))\, \mathrm{d}x\)}{\(\int _{a}^{b}(f(x) + g(x))\, \mathrm{d}x +\int _{ b}^{c}(f(x) - g(x))\, \mathrm{d}x\)}{}{}}
\LANGpl{
\Question{
Wskaż  wzór na pole obszaru zacieniowanego.}
{\(\int _{a}^{b}(f(x) - g(x))\, \mathrm{d}x +\int _{ b}^{c}(g(x) - f(x))\, \mathrm{d}x\)}{\(\int _{a}^{b}(g(x) - f(x))\, \mathrm{d}x +\int _{ b}^{c}(g(x) - f(x))\, \mathrm{d}x\)}{\(\int _{a}^{b}(f(x) - g(x))\, \mathrm{d}x +\int _{ b}^{c}(f(x) - g(x))\, \mathrm{d}x\)}{\(\int _{a}^{b}(f(x) + g(x))\, \mathrm{d}x +\int _{ b}^{c}(f(x) - g(x))\, \mathrm{d}x\)}{}{}}
\LANGcs{
\Question{
Vyjádřete obsah barevně vyznačené plochy omezené grafy funkcí
\(f\) a
\(g\) na
intervalu \(\langle a;c\rangle \).}
{\(\int _{a}^{b}(f(x) - g(x))\, \mathrm{d}x +\int _{ b}^{c}(g(x) - f(x))\, \mathrm{d}x\)}{\(\int _{a}^{b}(g(x) - f(x))\, \mathrm{d}x +\int _{ b}^{c}(g(x) - f(x))\, \mathrm{d}x\)}{\(\int _{a}^{b}(f(x) - g(x))\, \mathrm{d}x +\int _{ b}^{c}(f(x) - g(x))\, \mathrm{d}x\)}{\(\int _{a}^{b}(f(x) + g(x))\, \mathrm{d}x +\int _{ b}^{c}(f(x) - g(x))\, \mathrm{d}x\)}{}{}}
\LANGsk{
\Question{
Vyjadrite obsah farebne vyznačenej plochy vymedzenej grafmi funkcií
\(f\) a
\(g\) na
intervale \(\langle a;c\rangle \).}
{\(\int _{a}^{b}(f(x) - g(x))\, \mathrm{d}x +\int _{ b}^{c}(g(x) - f(x))\, \mathrm{d}x\)}{\(\int _{a}^{b}(g(x) - f(x))\, \mathrm{d}x +\int _{ b}^{c}(g(x) - f(x))\, \mathrm{d}x\)}{\(\int _{a}^{b}(f(x) - g(x))\, \mathrm{d}x +\int _{ b}^{c}(f(x) - g(x))\, \mathrm{d}x\)}{\(\int _{a}^{b}(f(x) + g(x))\, \mathrm{d}x +\int _{ b}^{c}(f(x) - g(x))\, \mathrm{d}x\)}{}{}}
\LANGes{
\Question{
Utilizando integrales escribe fórmula para el área de la región sombreada.}
{\(\int _{a}^{b}(f(x) - g(x))\, \mathrm{d}x +\int _{ b}^{c}(g(x) - f(x))\, \mathrm{d}x\)}{\(\int _{a}^{b}(g(x) - f(x))\, \mathrm{d}x +\int _{ b}^{c}(g(x) - f(x))\, \mathrm{d}x\)}{\(\int _{a}^{b}(f(x) - g(x))\, \mathrm{d}x +\int _{ b}^{c}(f(x) - g(x))\, \mathrm{d}x\)}{\(\int _{a}^{b}(f(x) + g(x))\, \mathrm{d}x +\int _{ b}^{c}(f(x) - g(x))\, \mathrm{d}x\)}{}{}}
\End

\Data\ID{9000065609}
\Updated{2020-07-15 03:07:37}
\Level{A}
\SubArea{Applications of definite integral}
\LANGen{
\Question{
Find the area of the region bounded by the curves
\(y = -x + 3\),
\(y = x^{2} - 3x\).}
{\(\frac{32}
 {3} \)}{\(8\)}{\(\frac{8}
{3}\)}{\(\frac{16}
 {3} \)}{}{}}
\LANGpl{
\Question{
Oblicz  pole obszaru ograniczonego krzywymi
\(y = -x + 3\),
\(y = x^{2} - 3x\).}
{\(\frac{32}
 {3} \)}{\(8\)}{\(\frac{8}
{3}\)}{\(\frac{16}
 {3} \)}{}{}}
\LANGcs{
\Question{
Vypočítejte obsah plochy ohraničené křivkami:
\(y = -x + 3\),
\(y = x^{2} - 3x\).}
{\(\frac{32}
 {3} \)}{\(8\)}{\(\frac{8}
{3}\)}{\(\frac{16}
 {3} \)}{}{}}
\LANGsk{
\Question{
Vypočítajte obsah plochy ohraničenej krivkami:
\(y = -x + 3\),
\(y = x^{2} - 3x\).}
{\(\frac{32}
 {3} \)}{\(8\)}{\(\frac{8}
{3}\)}{\(\frac{16}
 {3} \)}{}{}}
\LANGes{
\Question{
Encuentra el área de la superficie limitada por las curvas 
\(y = -x + 3\),
\(y = x^{2} - 3x\).}
{\(\frac{32}
 {3} \)}{\(8\)}{\(\frac{8}
{3}\)}{\(\frac{16}
 {3} \)}{}{}}
\End

\Data\ID{9000065610}
\Updated{2021-08-19 10:08:35}
\Level{A}
\SubArea{Applications of definite integral}
\LANGen{
\Question{
Using definite integral find the area of the triangle defined by the following three
inequalities 
\[ 
\begin{aligned}y& > 0, &
\\y& < x + 3,
\\y& < 3 - x.
\\ \end{aligned}
\]}
{\(\int _{-3}^{0}(x + 3)\, \mathrm{d}x +\int _{ 0}^{3}(3 - x)\, \mathrm{d}x\)}{\(\int _{0}^{3}(x + 3)\, \mathrm{d}x\)}{\(\int _{-3}^{3}(3 - x)\, \mathrm{d}x\)}{\(\int _{-3}^{0}(3 - x)\, \mathrm{d}x +\int _{ 0}^{3}(x + 3)\, \mathrm{d}x\)}{}{}}
\LANGpl{
\Question{
Oblicz całką  pole trójkąta określonego trzema podanymi nierównościami
\[ 
\begin{aligned}y& > 0, &
\\y& < x + 3,
\\y& < 3 - x.
\\ \end{aligned}
\]}
{\(\int _{-3}^{0}(x + 3)\, \mathrm{d}x +\int _{ 0}^{3}(3 - x)\, \mathrm{d}x\)}{\(\int _{0}^{3}(x + 3)\, \mathrm{d}x\)}{\(\int _{-3}^{3}(3 - x)\, \mathrm{d}x\)}{\(\int _{-3}^{0}(3 - x)\, \mathrm{d}x +\int _{ 0}^{3}(x + 3)\, \mathrm{d}x\)}{}{}}
\LANGcs{
\Question{
Vypočítejte pomocí určitého integrálu obsah trojúhelníku, který je popsaný
nerovnicemi: \(y > 0\),
\(y < x + 3\),
\(y < 3 - x\).}
{\(\int _{-3}^{0}(x + 3)\, \mathrm{d}x +\int _{ 0}^{3}(3 - x)\, \mathrm{d}x\)}{\(\int _{0}^{3}(x + 3)\, \mathrm{d}x\)}{\(\int _{-3}^{3}(3 - x)\, \mathrm{d}x\)}{\(\int _{-3}^{0}(3 - x)\, \mathrm{d}x +\int _{ 0}^{3}(x + 3)\, \mathrm{d}x\)}{}{}}
\LANGsk{
\Question{
Vypočítajte pomocou určitého integrálu obsah trojuholníka, ktorý je popísaný
nerovnicami: \(y > 0\),
\(y < x + 3\),
\(y < 3 - x\).}
{\(\int _{-3}^{0}(x + 3)\, \mathrm{d}x +\int _{ 0}^{3}(3 - x)\, \mathrm{d}x\)}{\(\int _{0}^{3}(x + 3)\, \mathrm{d}x\)}{\(\int _{-3}^{3}(3 - x)\, \mathrm{d}x\)}{\(\int _{-3}^{0}(3 - x)\, \mathrm{d}x +\int _{ 0}^{3}(x + 3)\, \mathrm{d}x\)}{}{}}
\LANGes{
\Question{
Utiizando la integral definida encuentra el área del triángulo definido por las siguientes tres desigualdades 
\[ 
\begin{aligned}y& > 0, &
\\y& < x + 3,
\\y& < 3 - x.
\\ \end{aligned}
\]}
{\(\int _{-3}^{0}(x + 3)\, \mathrm{d}x +\int _{ 0}^{3}(3 - x)\, \mathrm{d}x\)}{\(\int _{0}^{3}(x + 3)\, \mathrm{d}x\)}{\(\int _{-3}^{3}(3 - x)\, \mathrm{d}x\)}{\(\int _{-3}^{0}(3 - x)\, \mathrm{d}x +\int _{ 0}^{3}(x + 3)\, \mathrm{d}x\)}{}{}}
\End

\Data\ID{1003068201}
\Updated{2021-08-17 11:08:29}
\Level{B}
\SubArea{Applications of definite integral}
\LANGen{
\Question{
The value of the integral 
\[ \frac{4\pi}9\int\limits_0^3 x^2\mathrm{d}x \]
is a number which represents:}
{the volume of a cone with the base radius of \( 2\,\mathrm{cm} \) and the height of \( 3\,\mathrm{cm} \).}{the volume of a cone with the base radius of \( 3\,\mathrm{cm} \) and the height of \( 2\,\mathrm{cm} \).}{the volume of a sphere segment which is a part of a sphere with the radius of \( \frac23\,\mathrm{cm} \) and the height of \( 3\,\mathrm{cm} \).}{volume of a sphere segment which is a part of a sphere with the radius of \( 3\,\mathrm{cm} \) and the height of \( \frac23\,\mathrm{cm} \).}{}{}}
\LANGpl{
\Question{
Wartość całki
\[ \frac{4\pi}9\int\limits_0^3 x^2\mathrm{d}x \]
to liczba, która odpowiada:}
{objętości stożka, którego promień podstawy jest równy \( 2\,\mathrm{cm} \), a wysokość wynosi \( 3\,\mathrm{cm} \).}{objętości stożka, którego promień podstawy jest równy \( 3\,\mathrm{cm} \), a wysokość wynosi \( 2\,\mathrm{cm} \).}{objętości odcinka kuli, który jest częścią kuli, której  promień jest  równy \( \frac23\,\mathrm{cm} \), a wysokość wynosi \( 3\,\mathrm{cm} \).}{objętości odcinka kuli, który jest częścią kuli, której  promień jest  równy \( 3\,\mathrm{cm} \), a wysokość wynosi \( \frac23\,\mathrm{cm} \).}{}{}}
\LANGcs{
\Question{
Hodnotou výrazu
\[ \frac{4\pi}9\int\limits_0^3 x^2\mathrm{d}x \]
je číslo vyjadřující:}
{objem kužele o poloměru podstavy \( 2\,\mathrm{cm} \) a výšce \( 3\,\mathrm{cm} \).}{objem kužele o poloměru podstavy \( 3\,\mathrm{cm} \) a výšce  \( 2\,\mathrm{cm} \).}{objem kulové úseče, která je částí koule o poloměru \( \frac23\,\mathrm{cm} \) a má výšku \( 3\,\mathrm{cm} \).}{objem kulové úseče, která je částí koule o poloměru \( 3\,\mathrm{cm} \) a má výšku \( \frac23\,\mathrm{cm} \).}{}{}}
\LANGsk{
\Question{
Hodnotou výrazu
\[ \frac{4\pi}9\int\limits_0^3 x^2\mathrm{d}x \]
je číslo vyjadrujúce:}
{objem kužeľa s polomerom podstavy \( 2\,\mathrm{cm} \) a výškou \( 3\,\mathrm{cm} \).}{objem kužeľa s polomerom podstavy \( 3\,\mathrm{cm} \) a výškou  \( 2\,\mathrm{cm} \).}{objem guľovej úsečky, ktorá je časťou gule s polomerom \( \frac23\,\mathrm{cm} \) a má výšku \( 3\,\mathrm{cm} \).}{objem guľového úseku, ktorý je časťou gule s polomerom \( 3\,\mathrm{cm} \) a má výšku \( \frac23\,\mathrm{cm} \).}{}{}}
\LANGes{
\Question{
El valor de la integral  
\[ \frac{4\pi}9\int\limits_0^3 x^2\mathrm{d}x \]
es un número que representa:}
{el volumen de un cono cuya base tiene el radio de \( 2\,\mathrm{cm} \) y la altura son \( 3\,\mathrm{cm} \).}{el volumen de un cono cuya base tiene un radio de \( 3\,\mathrm{cm} \) y la altura son \( 2\,\mathrm{cm} \).}{el volumen de un casquete esférico que forma parte de una esfera cuyo radio es de \( \frac23\,\mathrm{cm} \) y la altura de \( 3\,\mathrm{cm} \).}{el volumen de un casquete esférico que forma parte de una esfera cuyo radio es de \( 3\,\mathrm{cm} \) y la altura de \( \frac23\,\mathrm{cm} \).}{}{}}
\End

\Data\ID{1003068202}
\Updated{2021-08-17 11:08:13}
\Level{B}
\SubArea{Applications of definite integral}
\LANGen{
\Question{
The value of the integral
\[ \pi\cdot\int\limits_0^6\left[9-(x-3)^2\right]\,\mathrm{d}x \]
is a number which represents:}
{the volume of a sphere with the radius of \( 3\,\mathrm{cm} \).}{the volume of a sphere with the radius of \( 6\,\mathrm{cm} \).}{the volume of a sphere with the diameter of \( 3\,\mathrm{cm} \).}{the volume of a semi-sphere with the radius of \( 3\,\mathrm{cm} \).}{}{}}
\LANGpl{
\Question{
Wartość całki
\[ \pi\cdot\int\limits_0^6\left[9-(x-3)^2\right]\,\mathrm{d}x \]
to liczba, która odpowiada:}
{objętości kuli o promieniu równym \( 3\,\mathrm{cm} \).}{objętości kuli o promieniu równym  \( 6\,\mathrm{cm} \).}{objętości kuli o średnicy równej \( 3\,\mathrm{cm} \).}{objętości półkuli o promieniu \( 3\,\mathrm{cm} \).}{}{}}
\LANGcs{
\Question{
Hodnotou výrazu 
\[ \pi\cdot\int\limits_0^6\left[9-(x-3)^2\right]\,\mathrm{d}x \]
je číslo vyjadřující:}
{objem koule o poloměru \( 3\,\mathrm{cm} \).}{objem koule o poloměru  \( 6\,\mathrm{cm} \).}{objem koule o průměru  \( 3\,\mathrm{cm} \).}{objem polokoule o poloměru  \( 3\,\mathrm{cm} \).}{}{}}
\LANGsk{
\Question{
Hodnotou výrazu 
\[ \pi\cdot\int\limits_0^6\left[9-(x-3)^2\right]\,\mathrm{d}x \]
je číslo vyjadrujúce:}
{objem gule s polomerom \( 3\,\mathrm{cm} \).}{objem gule s polomerom  \( 6\,\mathrm{cm} \).}{objem gule s priemerom  \( 3\,\mathrm{cm} \).}{objem pologule s polomerom  \( 3\,\mathrm{cm} \).}{}{}}
\LANGes{
\Question{
El valor de la integral 
\[ \pi\cdot\int\limits_0^6\left[9-(x-3)^2\right]\,\mathrm{d}x \]
es un número que representa:}
{el volumen de una esfera cuyo radio es de \( 3\,\mathrm{cm} \).}{el volumen de una esfera cuyo radio es de \( 6\,\mathrm{cm} \).}{el volumen de una esfera cuyo diámetro es de \( 3\,\mathrm{cm} \).}{el volumen de una semiesfera cuyo radio es de \( 3\,\mathrm{cm} \).}{}{}}
\End

\Data\ID{1003068203}
\Updated{2021-08-17 11:08:47}
\Level{B}
\SubArea{Applications of definite integral}
\LANGen{
\Question{
What is the volume of the solid that we get by rotating the curve
\[ y=\frac1{x^2} \]
about  the \( x \)-axis  on the interval \( [1;3] \)?}
{\( \frac{26}{81}\pi \)}{\( \frac{3^5-1}{3^6}\pi \)}{\( \frac{28}{81}\pi \)}{\( \frac{3^5+1}{3^6}\pi \)}{}{}}
\LANGpl{
\Question{
Wskaż objętość bryły otrzymanej w wyniku obrotu krzywej
\[ y=\frac1{x^2} \]
wokół osi \( x \) w przedziale  \( \langle1;3\rangle \)?}
{\( \frac{26}{81}\pi \)}{\( \frac{3^5-1}{3^6}\pi \)}{\( \frac{28}{81}\pi \)}{\( \frac{3^5+1}{3^6}\pi \)}{}{}}
\LANGcs{
\Question{
Jaký objem má těleso vzniklé rotací křivky
\[ y=\frac1{x^2} \]
kolem  \( x \)-ové osy  na intervalu\( \langle1;3\rangle \)?}
{\( \frac{26}{81}\pi \)}{\( \frac{3^5-1}{3^6}\pi \)}{\( \frac{28}{81}\pi \)}{\( \frac{3^5+1}{3^6}\pi \)}{}{}}
\LANGsk{
\Question{
Aký objem má teleso vzniknuté rotáciou krivky
\[ y=\frac1{x^2} \]
okolo osy  \( x \) na intervale\( \langle1;3\rangle \)?}
{\( \frac{26}{81}\pi \)}{\( \frac{3^5-1}{3^6}\pi \)}{\( \frac{28}{81}\pi \)}{\( \frac{3^5+1}{3^6}\pi \)}{}{}}
\LANGes{
\Question{
Calcula el volumen del sólido generado por la rotación de la curva
\[ y=\frac1{x^2} \]
alrededor del eje \( x \) en el intervalo \( [1;3] \).}
{\( \frac{26}{81}\pi \)}{\( \frac{3^5-1}{3^6}\pi \)}{\( \frac{28}{81}\pi \)}{\( \frac{3^5+1}{3^6}\pi \)}{}{}}
\End

\Data\ID{1003118702}
\Updated{2021-08-17 11:08:27}
\Level{B}
\SubArea{Applications of definite integral}
\LANGen{
\Question{
It is possible to calculate the volume of a sphere with the radius of \( 3 \) using a definite integral. Which of the following formulas is not correct?}
{\( \int\limits_{-3}^3\left(9-x^2\right)\,\mathrm{d}x \)}{\( \pi\int\limits_{-3}^3\left(9-x^2\right)\,\mathrm{d}x \)}{\( 2\pi\int\limits_{0}^3\left(9-x^2\right)\,\mathrm{d}x \)}{\( \pi\int\limits_{-3}^3\left(-\sqrt{9-x^2}\right)^2\,\mathrm{d}x \)}{}{}}
\LANGpl{
\Question{
Objętość kuli o promieniu  \( 3 \) można obliczyć za pomocą całki oznaczonej. Która z podanych formuł nie jest poprawna?}
{\( \int\limits_{-3}^3\left(9-x^2\right)\,\mathrm{d}x \)}{\( \pi\int\limits_{-3}^3\left(9-x^2\right)\,\mathrm{d}x \)}{\( 2\pi\int\limits_{0}^3\left(9-x^2\right)\,\mathrm{d}x \)}{\( \pi\int\limits_{-3}^3\left(-\sqrt{9-x^2}\right)^2\,\mathrm{d}x \)}{}{}}
\LANGcs{
\Question{
Pomocí určitého integrálu lze vypočítat objem koule o poloměru \( 3 \). Který z uvedených vzorců nelze použít?}
{\( \int\limits_{-3}^3\left(9-x^2\right)\,\mathrm{d}x \)}{\( \pi\int\limits_{-3}^3\left(9-x^2\right)\,\mathrm{d}x \)}{\( 2\pi\int\limits_{0}^3\left(9-x^2\right)\,\mathrm{d}x \)}{\( \pi\int\limits_{-3}^3\left(-\sqrt{9-x^2}\right)^2\,\mathrm{d}x \)}{}{}}
\LANGsk{
\Question{
Pomocou určitého integrálu sa dá vypočítať objem gule s polomerom \( 3 \). Ktorý z uvedených vzorcov sa nedá použiť?}
{\( \int\limits_{-3}^3\left(9-x^2\right)\,\mathrm{d}x \)}{\( \pi\int\limits_{-3}^3\left(9-x^2\right)\,\mathrm{d}x \)}{\( 2\pi\int\limits_{0}^3\left(9-x^2\right)\,\mathrm{d}x \)}{\( \pi\int\limits_{-3}^3\left(-\sqrt{9-x^2}\right)^2\,\mathrm{d}x \)}{}{}}
\LANGes{
\Question{
Es posible calcular el volumen de una esfera de radio de \( 3 \) usando una integral definida. ¿Cuál de las siguientes expresiones no es correcta?}
{\( \int\limits_{-3}^3\left(9-x^2\right)\,\mathrm{d}x \)}{\( \pi\int\limits_{-3}^3\left(9-x^2\right)\,\mathrm{d}x \)}{\( 2\pi\int\limits_{0}^3\left(9-x^2\right)\,\mathrm{d}x \)}{\( \pi\int\limits_{-3}^3\left(-\sqrt{9-x^2}\right)^2\,\mathrm{d}x \)}{}{}}
\End

\Data\ID{1003118703}
\Updated{2021-08-17 11:08:21}
\Level{B}
\SubArea{Applications of definite integral}
\LANGen{
\Question{
A right trapezoid is bounded by \( y=ax+1 \), \( x=0 \), \( x=6 \), and by the \( x \)-axis. Rotating the trapezoid around the \( x \)-axis we get a truncated cone. Find the value of the parameter \( a > 0 \) so that the volume of the truncated cone is \( 26\pi \).}
{\( a=\frac13 \)}{\( a=\frac12 \)}{\( a=3 \)}{\( a=2 \)}{}{}}
\LANGpl{
\Question{
Dany jest trapez prostokątny ograniczony przez  \( y=ax+1 \), \( x=0 \), \( x=6 \) oraz oś \( x \). Obracając trapez wokół osi \( x \) otrzymujemy stożek ścięty. Oblicz wartość parametru  \( a > 0 \) tak, aby objętość stożka ściętego wynosiła  \( 26\pi \).}
{\( a=\frac13 \)}{\( a=\frac12 \)}{\( a=3 \)}{\( a=2 \)}{}{}}
\LANGcs{
\Question{
Pravoúhlý lichoběžník je ohraničen přímkami \( y=ax+1 \), \( x=0 \), \( x=6 \) a osou \( x \). Jeho rotací kolem osy \( x \) vznikne komolý kužel. Určete hodnotu parametru \( a > 0 \) tak, aby byl objem tohoto komolého kužele \( 26\pi \).}
{\( a=\frac13 \)}{\( a=\frac12 \)}{\( a=3 \)}{\( a=2 \)}{}{}}
\LANGsk{
\Question{
Pravouhlý lichobežník je ohraničený priamkami \( y=ax+1 \), \( x=0 \), \( x=6 \) a osou \( x \). Jeho rotáciou okolo osy \( x \) vznikne zrezaný kužeľ. Určte hodnotu parametra \( a > 0 \) tak, aby bol objem tohoto zrezaného kužeľa \( 26\pi \).}
{\( a=\frac13 \)}{\( a=\frac12 \)}{\( a=3 \)}{\( a=2 \)}{}{}}
\LANGes{
\Question{
Un trapecio rectángulo está limitado por \( y=ax+1 \), \( x=0 \), \( x=6 \), y por el eje \( x \). Rotando el trapecio alrededor del eje \( x \) obtenemos un cono truncado. Encuentra el valor del parámetro  \( a > 0 \) para que el volumen del cono truncado sea \( 26\pi \).}
{\( a=\frac13 \)}{\( a=\frac12 \)}{\( a=3 \)}{\( a=2 \)}{}{}}
\End

\Data\ID{1003118705}
\Updated{2021-08-17 11:08:06}
\Level{B}
\SubArea{Applications of definite integral}
\LANGen{
\Question{
Peter and Jane both calculated the volume of a solid of revolution using a definite integral. Both chose a solid obtained by rotation of a line segment about the \( x \)-axis. The endpoints of Peter’s line segment are \( [0;1] \) and \( [5;4] \), the endpoints of Jane’s line segment are \( [0;3] \) and \( [5;0] \). Finally, they compared their calculated volumes. Which of the following statements is true?}
{Peter’s solid is \( 20\pi \) bigger.}{Jane’s solid is \( 20\pi \) bigger.}{Both solids have the same volume.}{The difference between Peter’s solid and Jane’s solid is \( 10\pi \).}{}{}}
\LANGpl{
\Question{
Piotr i Jan obliczyli objętość bryły obrotowej używając całki oznaczonej. Oboje wybrali bryłę uzyskaną przez obrót odcinka wokół osi \( x \). Punkty końcowe odcinka Piotra to  \( [0;1] \) oraz \( [5;4] \), punkty końcowe odcinka Jana to  \( [0;3] \) oraz  \( [5;0] \). Piotr i  Jan porównali swoje wyniki.  Wskaż zdanie prawdziwe.}
{Bryła Piotra jest większa o \( 20\pi \).}{Bryła Jana jest większa o \( 20\pi \).}{Bryły mają taką samą objętość.}{Różnica objętości pomiędzy bryłą Piotra i Jana wynosi \( 10\pi \).}{}{}}
\LANGcs{
\Question{
Petr a Jana počítali objem rotačního tělesa užitím určitého integrálu. Petr počítal objem tělesa vzniklého rotací úsečky s krajními body  \( [0;1] \) a \( [5;4] \) kolem osy \( x \). Jana počítala objem tělesa vzniklého rotací úsečky s krajními body  \( [0;3] \) a \( [5;0] \) rovněž kolem osy \( x \). Nemohli se pak dohodnout při porovnávání vypočítaných objemů. Které z uvedených tvrzení je pravdivé?}
{Petrovo těleso má objem o \( 20\pi \) větší.}{Janino těleso má objem o  \( 20\pi \) větší.}{Obě tělesa mají stejný objem.}{Rozdíl mezi objemy Petrova a Janina tělesa je \( 10\pi \).}{}{}}
\LANGsk{
\Question{
Peter a Jana počítali objem rotačného telesa použitím určitého integrálu. Peter počítal objem telesa vzniknutého rotáciou úsečky s krajnými bodmi  \( [0;1] \) a \( [5;4] \) okolo osy \( x \). Jana počítala objem telesa vzniknutého rotáciou úsečky s krajnými bodmi  \( [0;3] \) a \( [5;0] \) tiež okolo osy \( x \). Nemohli sa však dohodnúť pri porovnávaní vypočítaných objemov. Ktoré z uvedených tvrdení je pravdivé?}
{Petrovo teleso má objem o \( 20\pi \) väčší.}{Janine teleso má objem o  \( 20\pi \) väčší.}{Obidve telesa majú rovnaký objem.}{Rozdiel medzi objemami Petrovho a Janinho telesa je \( 10\pi \).}{}{}}
\LANGes{
\Question{
Pedro y Juana calcularon el volumen de un sólido de revolución mediante una integral definida. Ambos eligieron un sólido obtenido por rotación de un segmento alrededor del eje \( x \).  Pedro eligió el segmento de extremos \( [0;1] \) y \( [5;4] \),  y Juana eligió el de extremos \( [0;3] \) y \( [5;0] \). Finalmente compararon los volúmenes que cada uno había calculado.  ¿Cuál de los enunciados es verdadero?}
{El sólido de Pedro es \( 20\pi \) más grande.}{El sólido de Juana es \( 20\pi \) más grande.}{Ambos sólidos tienen el mismo volumen.}{La diferencia entre el sólido de Pedro y el de Juana es \( 10\pi \).}{}{}}
\End

\Data\ID{1003118706}
\Updated{2021-08-17 11:08:40}
\Level{B}
\SubArea{Applications of definite integral}
\LANGen{
\Question{
Consider a truncated cone with the diameters of its bases \( 2\,\mathrm{cm} \) and \( 10\,\mathrm{cm} \), and with the height \( 4\,\mathrm{cm} \).    Which of the following formulas cannot be used to calculate the volume of such truncated cone?}
{\( V=\pi\int\limits_0^4(5-x)\,\mathrm{d}x \)}{\( V=\pi\int\limits_0^4(5-x)^2\mathrm{d}x \)}{\( V=\frac{\pi}3\cdot4\cdot(25+5+1) \)}{\( V=\frac{\pi}3\cdot25\cdot5-\frac{\pi}3\cdot1\cdot1 \)}{}{}}
\LANGpl{
\Question{
Dany jest stożek ścięty, średnice jego podstaw są równe  \( 2\,\mathrm{cm} \) i \( 10\,\mathrm{cm} \), wysokość wynosi \( 4\,\mathrm{cm} \).  Które z poniższych wzorów nie mogą być użyte do obliczenia objętości tego stożka?}
{\( V=\pi\int\limits_0^4(5-x)\,\mathrm{d}x \)}{\( V=\pi\int\limits_0^4(5-x)^2\mathrm{d}x \)}{\( V=\frac{\pi}3\cdot4\cdot(25+5+1) \)}{\( V=\frac{\pi}3\cdot25\cdot5-\frac{\pi}3\cdot1\cdot1 \)}{}{}}
\LANGcs{
\Question{
Který z uvedených vzorců nelze použít pro výpočet objemu komolého kužele o výšce  \( 4\,\mathrm{cm} \), mají-li jeho podstavy průměry \( 2\,\mathrm{cm} \) a \( 10\,\mathrm{cm} \)?}
{\( V=\pi\int\limits_0^4(5-x)\,\mathrm{d}x \)}{\( V=\pi\int\limits_0^4(5-x)^2\mathrm{d}x \)}{\( V=\frac{\pi}3\cdot4\cdot(25+5+1) \)}{\( V=\frac{\pi}3\cdot25\cdot5-\frac{\pi}3\cdot1\cdot1 \)}{}{}}
\LANGsk{
\Question{
Ktorý z uvedených vzorcov sa nedá použiť pre výpočet objemu zrezaného kužeľa s výškou  \( 4\,\mathrm{cm} \), ak majú jeho podstavy priemery \( 2\,\mathrm{cm} \) a \( 10\,\mathrm{cm} \)?}
{\( V=\pi\int\limits_0^4(5-x)\,\mathrm{d}x \)}{\( V=\pi\int\limits_0^4(5-x)^2\mathrm{d}x \)}{\( V=\frac{\pi}3\cdot4\cdot(25+5+1) \)}{\( V=\frac{\pi}3\cdot25\cdot5-\frac{\pi}3\cdot1\cdot1 \)}{}{}}
\LANGes{
\Question{
Considera un cono truncado los diámetros de las bases de cual son \( 2\,\mathrm{cm} \) y \( 10\,\mathrm{cm} \), y cuya altura es \( 4\,\mathrm{cm} \). ¿Cuál de las fórmulas no debe ser usada para calcular el volumen de dicho cono truncado?}
{\( V=\pi\int\limits_0^4(5-x)\,\mathrm{d}x \)}{\( V=\pi\int\limits_0^4(5-x)^2\mathrm{d}x \)}{\( V=\frac{\pi}3\cdot4\cdot(25+5+1) \)}{\( V=\frac{\pi}3\cdot25\cdot5-\frac{\pi}3\cdot1\cdot1 \)}{}{}}
\End

\Data\ID{1103068301}
\Updated{2020-07-15 03:07:21}
\Level{B}
\SubArea{Applications of definite integral}
\IMG{1103068301_0.pdf}
\LANGen{
\Question{
Which of these formulas can be used to find the volume of the cone in the picture?}
{\( \pi\cdot\int\limits_0^4\left(-\frac14x+1\right)^2\,\mathrm{d}x \)}{\( \pi\cdot\int\limits_0^1\left(-4x+4\right)^2\,\mathrm{d}x \)}{\( 2\pi\cdot\int\limits_0^4\left(-\frac14x+1\right)^2\,\mathrm{d}x \)}{\( 2\pi\cdot\int\limits_0^1\left(-4x+4\right)^2\,\mathrm{d}x \)}{}{}}
\LANGpl{
\Question{
Które z tych wzorów można wykorzystać do obliczenia objętości stożka znajdującego się na rysunku?}
{\( \pi\cdot\int\limits_0^4\left(-\frac14x+1\right)^2\,\mathrm{d}x \)}{\( \pi\cdot\int\limits_0^1\left(-4x+4\right)^2\,\mathrm{d}x \)}{\( 2\pi\cdot\int\limits_0^4\left(-\frac14x+1\right)^2\,\mathrm{d}x \)}{\( 2\pi\cdot\int\limits_0^1\left(-4x+4\right)^2\,\mathrm{d}x \)}{}{}}
\LANGcs{
\Question{
Který vzorec lze použít pro výpočet objemu kužele z obrázku?}
{\( \pi\cdot\int\limits_0^4\left(-\frac14x+1\right)^2\,\mathrm{d}x \)}{\( \pi\cdot\int\limits_0^1\left(-4x+4\right)^2\,\mathrm{d}x \)}{\( 2\pi\cdot\int\limits_0^4\left(-\frac14x+1\right)^2\,\mathrm{d}x \)}{\( 2\pi\cdot\int\limits_0^1\left(-4x+4\right)^2\,\mathrm{d}x \)}{}{}}
\LANGsk{
\Question{
Ktorý vzorec sa dá použiť pre výpočet objemu kužeľa z obrázku?}
{\( \pi\cdot\int\limits_0^4\left(-\frac14x+1\right)^2\,\mathrm{d}x \)}{\( \pi\cdot\int\limits_0^1\left(-4x+4\right)^2\,\mathrm{d}x \)}{\( 2\pi\cdot\int\limits_0^4\left(-\frac14x+1\right)^2\,\mathrm{d}x \)}{\( 2\pi\cdot\int\limits_0^1\left(-4x+4\right)^2\,\mathrm{d}x \)}{}{}}
\LANGes{
\Question{
¿Cuál de las expresiones se puede usar para encontrar el volumen del cono en la imagen?}
{\( \pi\cdot\int\limits_0^4\left(-\frac14x+1\right)^2\,\mathrm{d}x \)}{\( \pi\cdot\int\limits_0^1\left(-4x+4\right)^2\,\mathrm{d}x \)}{\( 2\pi\cdot\int\limits_0^4\left(-\frac14x+1\right)^2\,\mathrm{d}x \)}{\( 2\pi\cdot\int\limits_0^1\left(-4x+4\right)^2\,\mathrm{d}x \)}{}{}}
\End

\Data\ID{1103068302}
\Updated{2021-08-15 09:08:54}
\Level{B}
\SubArea{Applications of definite integral}
\IMG{1103068302_0.pdf}
\LANGen{
\Question{
Which of these formulas can be used to find the volume of the cylinder in the picture? Points \( [0; 0; 0] \) and \( [4;0;0] \) are located in the centers of the cylinder bases.}
{\( \pi\cdot\int\limits_0^43^2\,\mathrm{d}x \)}{\( \pi\cdot\int\limits_0^34^2\,\mathrm{d}x \)}{\( \pi\cdot\int\limits_0^43\,\mathrm{d}x \)}{\( \pi\cdot\int\limits_{-4}^49\,\mathrm{d}x \)}{}{}}
\LANGpl{
\Question{
Które z tych wzorów można wykorzystać do obliczenia objętości walca znajdującego się na rysunku? Punkty  \( [0; 0; 0] \) i \( [4;0;0] \)  znajdują się w środku podstaw walca.}
{\( \pi\cdot\int\limits_0^43^2\,\mathrm{d}x \)}{\( \pi\cdot\int\limits_0^34^2\,\mathrm{d}x \)}{\( \pi\cdot\int\limits_0^43\,\mathrm{d}x \)}{\( \pi\cdot\int\limits_{-4}^49\,\mathrm{d}x \)}{}{}}
\LANGcs{
\Question{
Který vzorec lze použít pro výpočet objemu válce z obrázku? Body \( [0; 0; 0] \) a \( [4;0;0] \) jsou středy podstav.}
{\( \pi\cdot\int\limits_0^43^2\,\mathrm{d}x \)}{\( \pi\cdot\int\limits_0^34^2\,\mathrm{d}x \)}{\( \pi\cdot\int\limits_0^43\,\mathrm{d}x \)}{\( \pi\cdot\int\limits_{-4}^49\,\mathrm{d}x \)}{}{}}
\LANGsk{
\Question{
Ktorý vzorec sa dá použiť pre výpočet objemu valca z obrázku? Body \( [0; 0; 0] \) a \( [4;0;0] \) sú stredy podstáv.}
{\( \pi\cdot\int\limits_0^43^2\,\mathrm{d}x \)}{\( \pi\cdot\int\limits_0^34^2\,\mathrm{d}x \)}{\( \pi\cdot\int\limits_0^43\,\mathrm{d}x \)}{\( \pi\cdot\int\limits_{-4}^49\,\mathrm{d}x \)}{}{}}
\LANGes{
\Question{
¿Cuál de las fórmulas se puede usar para encontrar el volumen del cilindro de la imagen? Los puntos \( (0; 0; 0) \) y \( (4;0;0) \) están situados en los centros de las bases del cilindro.}
{\( \pi\cdot\int\limits_0^43^2\,\mathrm{d}x \)}{\( \pi\cdot\int\limits_0^34^2\,\mathrm{d}x \)}{\( \pi\cdot\int\limits_0^43\,\mathrm{d}x \)}{\( \pi\cdot\int\limits_{-4}^49\,\mathrm{d}x \)}{}{}}
\End

\Data\ID{1103068303}
\Updated{2021-08-15 09:08:27}
\Level{B}
\SubArea{Applications of definite integral}
\IMG{1103068303_0.pdf}
\LANGen{
\Question{
Which  of the following  formulas will not give the volume of the solid created by the rotation of the red area about  the \( x \)-axis (see the picture)?}
{\( \pi\cdot\int\limits_{\frac{\pi}4}^{\frac{3\pi}4}\sin x\,\mathrm{d}x \)}{\( \pi\cdot\int\limits_{\frac{\pi}4}^{\frac{3\pi}4}\sin^2x\,\mathrm{d}x \)}{\( 2\pi\cdot\int\limits_{\frac{\pi}2}^{\frac{3\pi}4}\sin^2x\,\mathrm{d}x \)}{\( \pi\cdot\int\limits_{\frac{9\pi}4}^{\frac{11\pi}4}\sin^2x\,\mathrm{d}x \)}{}{}}
\LANGpl{
\Question{
Który z poniższych wzorów nie można wykorzystać do obliczenia objętości bryły uzyskanej przez obrót powierzchni zaznaczonej kolorem czerwonym wokół osi \( x \) (spójrz na rysunek)?}
{\( \pi\cdot\int\limits_{\frac{\pi}4}^{\frac{3\pi}4}\sin x\,\mathrm{d}x \)}{\( \pi\cdot\int\limits_{\frac{\pi}4}^{\frac{3\pi}4}\sin^2x\,\mathrm{d}x \)}{\( 2\pi\cdot\int\limits_{\frac{\pi}2}^{\frac{3\pi}4}\sin^2x\,\mathrm{d}x \)}{\( \pi\cdot\int\limits_{\frac{9\pi}4}^{\frac{11\pi}4}\sin^2x\,\mathrm{d}x \)}{}{}}
\LANGcs{
\Question{
Jaký objem nebude mít těleso vzniklé rotací vyznačeného červeného útvaru kolem osy \( x \)?}
{\( \pi\cdot\int\limits_{\frac{\pi}4}^{\frac{3\pi}4}\sin x\,\mathrm{d}x \)}{\( \pi\cdot\int\limits_{\frac{\pi}4}^{\frac{3\pi}4}\sin^2x\,\mathrm{d}x \)}{\( 2\pi\cdot\int\limits_{\frac{\pi}2}^{\frac{3\pi}4}\sin^2x\,\mathrm{d}x \)}{\( \pi\cdot\int\limits_{\frac{9\pi}4}^{\frac{11\pi}4}\sin^2x\,\mathrm{d}x \)}{}{}}
\LANGsk{
\Question{
Aký objem nebude mať teleso vzniknuté rotáciou vyznačeného červeného útvaru okolo osy \( x \)?}
{\( \pi\cdot\int\limits_{\frac{\pi}4}^{\frac{3\pi}4}\sin x\,\mathrm{d}x \)}{\( \pi\cdot\int\limits_{\frac{\pi}4}^{\frac{3\pi}4}\sin^2x\,\mathrm{d}x \)}{\( 2\pi\cdot\int\limits_{\frac{\pi}2}^{\frac{3\pi}4}\sin^2x\,\mathrm{d}x \)}{\( \pi\cdot\int\limits_{\frac{9\pi}4}^{\frac{11\pi}4}\sin^2x\,\mathrm{d}x \)}{}{}}
\LANGes{
\Question{
¿Cuál de las siguientes fórmulas no va a dar el volumen del sólido generado por rotación del área roja alrededor del eje \( x \) (mira la imagen)?}
{\( \pi\cdot\int\limits_{\frac{\pi}4}^{\frac{3\pi}4}\sin x\,\mathrm{d}x \)}{\( \pi\cdot\int\limits_{\frac{\pi}4}^{\frac{3\pi}4}\sin^2x\,\mathrm{d}x \)}{\( 2\pi\cdot\int\limits_{\frac{\pi}2}^{\frac{3\pi}4}\sin^2x\,\mathrm{d}x \)}{\( \pi\cdot\int\limits_{\frac{9\pi}4}^{\frac{11\pi}4}\sin^2x\,\mathrm{d}x \)}{}{}}
\End

\Data\ID{1103068304}
\Updated{2020-07-15 03:07:21}
\Level{B}
\SubArea{Applications of definite integral}
\IMG{1103068304_0.pdf}
\LANGen{
\Question{
Find the volume of the solid obtained  by the rotation of the yellow area about  the \( x \)-axis (see the picture).}
{\( \frac56\pi \)}{\( \frac76\pi \)}{\( \pi\cdot\ln 6 \)}{\( \frac{35}{36}\pi \)}{}{}}
\LANGpl{
\Question{
Oblicz objętość bryły powstałej z obrotu obszaru zaznaczonego żółtym kolorem wokół osi  \( x \). (Zobacz na rysunek.)}
{\( \frac56\pi \)}{\( \frac76\pi \)}{\( \pi\cdot\ln 6 \)}{\( \frac{35}{36}\pi \)}{}{}}
\LANGcs{
\Question{
Určete objem tělesa vzniklého rotací vyznačeného žlutého útvaru kolem osy \( x \).}
{\( \frac56\pi \)}{\( \frac76\pi \)}{\( \pi\cdot\ln 6 \)}{\( \frac{35}{36}\pi \)}{}{}}
\LANGsk{
\Question{
Určte objem telesa vzniknutého rotáciou vyznačeného žltého útvaru okolo osy \( x \).}
{\( \frac56\pi \)}{\( \frac76\pi \)}{\( \pi\cdot\ln 6 \)}{\( \frac{35}{36}\pi \)}{}{}}
\LANGes{
\Question{
Encuentra el volumen del sólido obtenido por rotación del área amarilla alrededor del eje \( x \) (mira la imagen).}
{\( \frac56\pi \)}{\( \frac76\pi \)}{\( \pi\cdot\ln 6 \)}{\( \frac{35}{36}\pi \)}{}{}}
\End

\Data\ID{1103118701}
\Updated{2021-08-15 09:08:03}
\Level{B}
\SubArea{Applications of definite integral}
\IMG{1103118701.pdf}
\LANGen{
\Question{
Find the volume of the truncated cone obtained by the rotation of the shaded trapezium (see the picture) about the \( x \)-axis.}
{\( 42\pi \)}{\( 42 \)}{\( 16\pi \)}{\( 16 \)}{}{}}
\LANGpl{
\Question{
Oblicz objętość ściętego stożka uzyskanego przez obrót zacienionego trapezu (patrz obrazek) wokół osi  \( x \).}
{\( 42\pi \)}{\( 42 \)}{\( 16\pi \)}{\( 16 \)}{}{}}
\LANGcs{
\Question{
Jaký objem bude mít komolý kužel, který vznikne rotací vyšrafovaného lichoběžníku kolem osy \( x \)?}
{\( 42\pi \)}{\( 42 \)}{\( 16\pi \)}{\( 16 \)}{}{}}
\LANGsk{
\Question{
Aký objem bude mať zrezaný kužeľ, ktorý vznikne rotáciou vyšrafovaného lichobežníka okolo osy \( x \)?}
{\( 42\pi \)}{\( 42 \)}{\( 16\pi \)}{\( 16 \)}{}{}}
\LANGes{
\Question{
Encuentra el volumen del cono truncado obtenido por rotación del trapecio marcado (mira la imagen) alrededor del eje \( x \).}
{\( 42\pi \)}{\( 42 \)}{\( 16\pi \)}{\( 16 \)}{}{}}
\End

\Data\ID{1103118704}
\Updated{2021-08-19 10:08:51}
\Level{B}
\SubArea{Applications of definite integral}
\IMG{1103118704_0.pdf}
\LANGen{
\Question{
Which of the given equations defines the line that together with \( x=0 \) and \( x \)-axis bounds the right triangle, if by rotation of this triangle about the \( x \)-axis the cone of height \( 10 \) is obtained as indicated in the picture?}
{\( y=-0.4x+4 \)}{\( y=-2.5x+10 \)}{\( y=4x+10 \)}{\( y=10x+4 \)}{}{}}
\LANGpl{
\Question{
Które z podanych równań określa prostą, która wraz z  \( x=0 \) i osią \( x \) ogranicza trójkąt prostokątny,  jeśli w wyniku obrotu tego trójkąta wokół osi \( x \) powstaje stożek o wysokości \( 10 \)?}
{\( y=-0{,}4x+4 \)}{\( y=-2{,}5x+10 \)}{\( y=4x+10 \)}{\( y=10x+4 \)}{}{}}
\LANGcs{
\Question{
Kterou z uvedených rovnic je dána přímka vymezující s přímkou \( x=0 \) a osou \( x \) pravoúhlý trojúhelník, jehož rotací kolem osy \( x \) vzniká zakreslený kužel výšky \( 10 \)?}
{\( y=-0{,}4x+4 \)}{\( y=-2{,}5x+10 \)}{\( y=4x+10 \)}{\( y=10x+4 \)}{}{}}
\LANGsk{
\Question{
Ktorou z uvedených rovníc je daná priamka vymedzujúca s priamkou \( x=0 \) a osou \( x \) pravouhlý trojuholník, ktorého rotáciou okolo osy \( x \) vzniká zakreslený kužeľ výšky \( 10 \)?}
{\( y=-0{,}4x+4 \)}{\( y=-2{,}5x+10 \)}{\( y=4x+10 \)}{\( y=10x+4 \)}{}{}}
\LANGes{
\Question{
¿Cuál de las ecuaciones dadas define la recta que junto con \( x=0 \) y el eje \( x \) limita el triángulo rectángulo? Sabemos que  por rotación de este triángulo alrededor del eje \( x \) se obtiene un cono de altura \( 10 \) como esta indicado en la imagen.}
{\( y=-0.4x+4 \)}{\( y=-2.5x+10 \)}{\( y=4x+10 \)}{\( y=10x+4 \)}{}{}}
\End

\Data\ID{1103124403}
\Updated{2020-07-22 06:07:59}
\Level{B}
\SubArea{Applications of definite integral}
\IMG{1103124403.pdf}
\LANGen{
\Question{
Let a solid be obtained by rotating the yellow region about the $x$-axis (see the picture). Find its volume.}
{$85.5\pi$}{$85.5$}{$27$}{$27\pi$}{}{}}
\LANGpl{
\Question{
Dana jest bryła uzyskana przez obrót obszaru zaznaczonego żółtym kolorem wokół osi  $x$ (spójrz na rysunek). Oblicz jej objętość.}
{$85{,}5\pi$}{$85{,}5$}{$27$}{$27\pi$}{}{}}
\LANGcs{
\Question{
Jaký objem má těleso, které vznikne rotací žlutě vyznačeného obrazce kolem osy $x$?}
{$85{,}5\pi$}{$85{,}5$}{$27$}{$27\pi$}{}{}}
\LANGsk{
\Question{
Aký objem má teleso, ktoré vznikne rotáciou žlto vyznačeného obrazca okolo osi $x$?}
{$85{,}5\pi$}{$85{,}5$}{$27$}{$27\pi$}{}{}}
\LANGes{
\Question{
Encuentra el volumen del sólido obtenido por rotación de la región amarilla alrededor del eje $x$ (mira la imagen).}
{$85.5\pi$}{$85.5$}{$27$}{$27\pi$}{}{}}
\End

\Data\ID{9000100001}
\Updated{2020-08-03 04:08:23}
\Level{B}
\SubArea{Applications of definite integral}
\IMG{001000_01-figure0.pdf}
\LANGen{
\Question{
The function \(f(x) = 3 - 2x\) 
is graphed in the picture. Consider the region between the graph of the function on the
interval \([ 0;\, 1.5] \) 
and the axes. Determine the solid of revolution obtained by revolving this region about
\(y\)-axis}
{A cone with the base of radius \(1.5\).}{A cone with the base of radius \(3\).}{A pyramid of the height \(1.5\).}{A pyramid of the height \(3\).}{}{}}
\LANGpl{
\Question{
Wykres przedstawia funkcję  \(f\colon y = 3 - 2x\).  Jaka bryła obrotowa powstanie z obrotu obszaru ograniczonego przez wykres danej  funkcji na przedziale 
 \([ 0;\, 1.5] \) 
oraz  osiami x i y wokół osi 
\(y\)?}
{Stożek o promieniu podstawy  \(1.5\).}{Stożek o promieniu podstawy \(3\).}{Ostrosłup o wysokości \(1.5\).}{Ostrosłup o wysokości \(3\).}{}{}}
\LANGcs{
\Question{
Na obrázku je graf funkce \(f\colon y = 3 - 2x\).
Jaké těleso vznikne rotací rovinného obrazce ohraničeného osou
\(x\), osou
\(y\) a grafem
funkce \(f\) na
intervalu \(\langle 0;\, 1{,}5\rangle \) 
kolem osy \(y\)?}
{Kužel s poloměrem podstavy \(1{,}5\).}{Kužel s poloměrem podstavy \(3\).}{Jehlan s tělesovou výškou \(1{,}5\).}{Jehlan s tělesovou výškou \(3\).}{}{}}
\LANGsk{
\Question{
Na obrázku je graf funkcie \(f\colon y = 3 - 2x\).
Aké teleso vznikne rotáciou rovinného obrazca ohraničeného osou
\(x\), osou
\(y\) a grafom
funkcie \(f\) na
intervale \(\langle 0;\, 1{,}5\rangle \) 
okolo osy \(y\)?}
{Kužeľ s polomerom podstavy \(1{,}5\).}{Kužeľ s polomerom podstavy \(3\).}{Ihlan s telesovou výškou \(1{,}5\).}{Ihlan s telesovou výškou \(3\).}{}{}}
\LANGes{
\Question{
En la imagen se puede ver la gráfica de la función \(f(x) = 3 - 2x\). Consideremos la región entre la gráfica de la función en el intervalo \([ 0;\, 1.5] \) 
y los ejes. Determina el sólido de revolución obtenido por rotación de esta región alrededor del eje \(y\).}
{Cono con el radio de la base igual a \(1.5\).}{Cono con el radio de la base igual a \(3\).}{Pirámide de la altura igual a \(1.5\).}{Pirámide de la altura igual a \(3\).}{}{}}
\End

\Data\ID{9000100002}
\Updated{2020-09-23 12:09:54}
\Level{B}
\SubArea{Applications of definite integral}
\IMG{001000_02-figure0.pdf}
\LANGen{
\Question{
The function \(f(x) = 3 - 2x\) 
is graphed in the picture. Consider the region between the graph of the function \(f\), the
\(x\)-axis and
the lines \(x = 1\) 
and \(x = -1\).
Find the volume of the solid of revolution obtained by revolving this region about
\(x\)-axis.}
{\(\frac{62}
 {3} \pi \)}{\(6\pi \)}{\(12\pi \)}{\(\frac{8}
{3}\pi \)}{}{}}
\LANGpl{
\Question{
Wykres przedstawia funkcję \(f\colon y = 3 - 2x\) 
Jaka jest objętość  bryły obrotowej, która  powstanie z obrotu obszaru ograniczonego przez wykres danej funkcji na przedziale 
 \([ 0;\, 1.5] \), osią 
\(x\) oraz 
prostymi  \(x = 1\) 
i \(x = -1\) wokół osi 
\(x\).}
{\(\frac{62}
 {3} \pi \)}{\(6\pi \)}{\(12\pi \)}{\(\frac{8}
{3}\pi \)}{}{}}
\LANGcs{
\Question{
Na obrázku je graf funkce \(f\colon y = 3 - 2x\).
Jaký je objem tělesa, které vznikne rotací rovinného obrazce ohraničeného osou
\(x\), grafem
funkce \(f\) a
přímkami \(x = -1\) 
a \(x = 1\) kolem
osy \(x\)?}
{\(\frac{62}
 {3} \pi \)}{\(6\pi \)}{\(12\pi \)}{\(\frac{8}
{3}\pi \)}{}{}}
\LANGsk{
\Question{
Na obrázku je graf funkcie \(f\colon y = 3 - 2x\).
Aký je objem telesa, ktoré vznikne rotáciou rovinného obrazca ohraničeného osou
\(x\), grafom
funkcie \(f\) a
priamkami \(x = -1\) 
a \(x = 1\) okolo
osy \(x\)?}
{\(\frac{62}
 {3} \pi \)}{\(6\pi \)}{\(12\pi \)}{\(\frac{8}
{3}\pi \)}{}{}}
\LANGes{
\Question{
En la imagen se puede ver la gráfica de la función \(f(x) = 3 - 2x\). Consideremos la región entre la gráfica de la función \(f\), el eje \(x\) y
las rectas \(x = 1\) 
y \(x = -1\).
Halla el volumen del sólido de revolución obtenido por rotación de esta región alrededor del eje
\(x\).}
{\(\frac{62}
 {3} \pi \)}{\(6\pi \)}{\(12\pi \)}{\(\frac{8}
{3}\pi \)}{}{}}
\End

\Data\ID{9000100003}
\Updated{2020-08-03 04:08:08}
\Level{B}
\SubArea{Applications of definite integral}
\IMG{001000_03-figure0.pdf}
\LANGen{
\Question{
The function \(f(x) = x^{2} + 2\) 
is graphed in the picture. Consider the region between the graph of the function on the interval
\([ 0;\, 1] \), both axes
and the line \(x = 1\).
Find the formula for the volume of the solid of revolution obtained by revolving this region
about \(y\)-axis.}
{\(V =\pi \int _{ 0}^{3}1\, \mathrm{d}y -\pi \int _{2}^{3}(\sqrt{y - 2})^{2}\, \mathrm{d}y\)}{\(V =\pi \int _{ 0}^{3}(\sqrt{y - 2})^{2}\, \mathrm{d}y\)}{\(V =\pi \int _{ 2}^{3}(\sqrt{y - 2})^{2}\, \mathrm{d}y -\pi \int _{0}^{3}1\, \mathrm{d}y\)}{\(V =\pi \int _{ 2}^{3}(\sqrt{y - 2})^{2}\, \mathrm{d}y\)}{}{}}
\LANGpl{
\Question{
Wykres przedstawia funkcję \(f\colon y = x^{2} + 2\). 
Wskaż wzór na objętość bryły obrotowej, która  powstanie z obrotu obszaru ograniczonego przez wykres danej funkcji na przedziale
\([ 0;\, 1] \), obiema osiami
oraz prostą  \(x = 1\) wokół osi
\(y\).}
{\(V =\pi \int _{ 0}^{3}1\, \mathrm{d}y -\pi \int _{2}^{3}(\sqrt{y - 2})^{2}\, \mathrm{d}y\)}{\(V =\pi \int _{ 0}^{3}(\sqrt{y - 2})^{2}\, \mathrm{d}y\)}{\(V =\pi \int _{ 2}^{3}(\sqrt{y - 2})^{2}\, \mathrm{d}y -\pi \int _{0}^{3}1\, \mathrm{d}y\)}{\(V =\pi \int _{ 2}^{3}(\sqrt{y - 2})^{2}\, \mathrm{d}y\)}{}{}}
\LANGcs{
\Question{
Na obrázku je graf funkce \(f\colon y = x^{2} + 2\).
Pro objem tělesa, které vznikne rotací rovinného obrazce ohraničeného osou
\(x\), osou
\(y\), grafem
funkce \(f\) na
intervalu \(\langle 0;\, 1\rangle \) a
přímkou \(x = 1\) 
kolem osy \(y\) 
platí vztah:}
{\(V =\pi \int _{ 0}^{3}1\, \mathrm{d}y -\pi \int _{2}^{3}(\sqrt{y - 2})^{2}\, \mathrm{d}y\)}{\(V =\pi \int _{ 0}^{3}(\sqrt{y - 2})^{2}\, \mathrm{d}y\)}{\(V =\pi \int _{ 2}^{3}(\sqrt{y - 2})^{2}\, \mathrm{d}y -\pi \int _{0}^{3}1\, \mathrm{d}y\)}{\(V =\pi \int _{ 2}^{3}(\sqrt{y - 2})^{2}\, \mathrm{d}y\)}{}{}}
\LANGsk{
\Question{
Na obrázku je graf funkcie \(f\colon y = x^{2} + 2\).
Pre objem telesa, ktoré vznikne rotáciou rovinného obrazca ohraničeného osou
\(x\), osou
\(y\), grafom
funkcie \(f\) na
intervale \(\langle 0;\, 1\rangle \) a
priamkou \(x = 1\) 
okolo osy \(y\) 
platí vzťah:}
{\(V =\pi \int _{ 0}^{3}1\, \mathrm{d}y -\pi \int _{2}^{3}(\sqrt{y - 2})^{2}\, \mathrm{d}y\)}{\(V =\pi \int _{ 0}^{3}(\sqrt{y - 2})^{2}\, \mathrm{d}y\)}{\(V =\pi \int _{ 2}^{3}(\sqrt{y - 2})^{2}\, \mathrm{d}y -\pi \int _{0}^{3}1\, \mathrm{d}y\)}{\(V =\pi \int _{ 2}^{3}(\sqrt{y - 2})^{2}\, \mathrm{d}y\)}{}{}}
\LANGes{
\Question{
En la imagen se puede ver la gráfica de la función \(f(x) = x^{2} + 2\). Consideremos la región entre la gráfica de la función en el intervalo \([ 0;\, 1] \), los dos ejes y la recta \(x = 1\).
Busca la fórmula del volumen del sólido de revolución obtenido por rotación de esta región alrededor del eje \(y\).}
{\(V =\pi \int _{ 0}^{3}1\, \mathrm{d}y -\pi \int _{2}^{3}(\sqrt{y - 2})^{2}\, \mathrm{d}y\)}{\(V =\pi \int _{ 0}^{3}(\sqrt{y - 2})^{2}\, \mathrm{d}y\)}{\(V =\pi \int _{ 2}^{3}(\sqrt{y - 2})^{2}\, \mathrm{d}y -\pi \int _{0}^{3}1\, \mathrm{d}y\)}{\(V =\pi \int _{ 2}^{3}(\sqrt{y - 2})^{2}\, \mathrm{d}y\)}{}{}}
\End

\Data\ID{9000100004}
\Updated{2021-08-19 10:08:51}
\Level{B}
\SubArea{Applications of definite integral}
\IMG{001000_04-figure0.pdf}
\LANGen{
\Question{
The function \(f(x)= x^{2} + 2\) 
is graphed in the picture. Consider the region bounded by the graph of the function, both axes
and the line \(x = -1\).
Determine the solid of revolution obtained by revolving this region about
\(x\)-axis.}
{A general solid which is neither cone nor cylinder.}{Cone with base of radius \(1\).}{Cylinder with base of radius \(2\).}{Cone with base of radius \(2\).}{}{}}
\LANGpl{
\Question{
Wykres przedstawia funkcję \(f\colon y = x^{2} + 2\). 
Jaka bryła obrotowa powstanie z obrotu obszaru ograniczonego przez wykres danej funkcji, obiema osiami
oraz prostą \(x = -1\) wokół osi 
\(x\).}
{Bryła jednorodna nie będąca, ani stożkiem, ani walcem.}{Stożek o promieniu podstawy \(1\).}{Walec  o promieniu podstawy \(2\).}{Stożek o promieniu podstawy \(2\).}{}{}}
\LANGcs{
\Question{
Na obrázku je graf funkce \(f\colon y = x^{2} + 2\).
Jaké těleso vznikne rotací rovinného obrazce ohraničeného osou
\(x\), osou
\(y\), grafem
funkce \(f\) a
přímkou \(x = -1\) 
kolem osy \(x\)?}
{Těleso různé od kužele a válce.}{Kužel s poloměrem podstavy \(1\).}{Válec s poloměrem podstavy \(2\).}{Kužel s poloměrem podstavy \(2\).}{}{}}
\LANGsk{
\Question{
Na obrázku je graf funkcie \(f\colon y = x^{2} + 2\).
Aké teleso vznikne rotáciou rovinného obrazca ohraničeného osou
\(x\), osou
\(y\), grafom
funkcie \(f\) a
priamkou \(x = -1\) 
okolo osy \(x\)?}
{Teleso rôzne od kužeľa a valca.}{Kužeľ s polomerom podstavy \(1\).}{Valec s polomerom podstavy \(2\).}{Kužeľ s polomerom podstavy \(2\).}{}{}}
\LANGes{
\Question{
En la imagen se puede ver la gráfica de la función \(f(x)= x^{2} + 2\). Consideremos la región limitada por la gráfica de la función, los dos ejes y la recta \(x = -1\).
Determina el sólido de revolución obtenido por rotación de esta región alrededor del eje \(x\).}
{Un sólido general que no es ni cono ni cilindro.}{Cono con el radio de la base igual a \(1\).}{Cilindro con el radio de la base igual a \(2\).}{Cono con el radio de la base igual a \(2\).}{}{}}
\End

\Data\ID{9000100005}
\Updated{2020-08-03 04:08:43}
\Level{B}
\SubArea{Applications of definite integral}
\IMG{001000_05-figure0.pdf}
\LANGen{
\Question{
The function \(f(x) = 1\) 
is graphed in the picture. Determine the solid of revolution with volume given by the
following formula. 
\[ 
 \pi \int _{-1}^{1}f^{2}(x)\, \mathrm{d}x
\]}
{Cylinder of base radius \(1\) 
and height \(2\).}{Cone of base radius \(1\) 
and height \(2\).}{Cone of base radius \(2\) 
and height \(1\).}{Cylinder of base radius \(2\) 
and height \(1\).}{}{}}
\LANGpl{
\Question{
Wykres przedstawia funkcję \(f\colon y = 1\). 
Wskaż bryłę obrotową o objętości wyrażonej wzorem.
\[ 
 \pi \int _{-1}^{1}f^{2}(x)\, \mathrm{d}x
\]}
{Walec  o promieniu podstawy \(1\) 
i wysokości  \(2\).}{Stożek o promieniu podstawy \(1\) 
i wysokości  \(2\).}{Stożek o promieniu podstawy \(2\) 
i wysokości \(1\).}{Walec  o promieniu podstawy \(2\) 
i wysokości \(1\).}{}{}}
\LANGcs{
\Question{
Na obrázku je graf funkce \(f\colon y = 1\).
Určete těleso, jehož objem vypočítáme vztahem
\(\pi \int _{-1}^{1}f^{2}(x)\, \mathrm{d}x\).}
{Válec o poloměru podstavy \(1\) 
a výšce \(2\).}{Kužel o poloměru podstavy \(1\) 
a výšce \(2\).}{Kužel o poloměru podstavy \(2\) 
a výšce \(1\).}{Válec o poloměru podstavy \(2\) 
a výšce \(1\).}{}{}}
\LANGsk{
\Question{
Na obrázku je graf funkcie \(f\colon y = 1\).
Určte teleso, ktorého objem vypočítame vzťahom
\(\pi \int _{-1}^{1}f^{2}(x)\, \mathrm{d}x\).}
{Valec s polomerom podstavy \(1\) 
a výškou \(2\).}{Kužeľ s polomerom podstavy \(1\) 
a výškou \(2\).}{Kužeľ o polomerom podstavy \(2\) 
a výškou \(1\).}{Valec s polomerom podstavy \(2\) 
a výškou \(1\).}{}{}}
\LANGes{
\Question{
En la imagen se puede ver la gráfica de la función \(f(x) = 1\). Determina el sólido de revolución con el volumen dado por esta fórmula. 
\[ 
 \pi \int _{-1}^{1}f^{2}(x)\, \mathrm{d}x
\]}
{Cilindro del radio de la base igual a \(1\) 
y la altura a \(2\).}{Cono del radio de la base igual a \(1\) 
y la altura a \(2\).}{Cono del radio de la base igual a \(2\) 
y la altura a \(1\).}{Cilindro del radio de la base igual a \(2\) 
y la altura a \(1\).}{}{}}
\End

\Data\ID{9000100006}
\Updated{2021-08-19 10:08:17}
\Level{B}
\SubArea{Applications of definite integral}
\IMG{001000_06-figure0.pdf}
\LANGen{
\Question{
The function \(f(x)= \sqrt{x}\) 
is graphed in the picture. Consider the region bounded by the graph of
\(f\) on
\([ 1;\, 4] \), lines
\(x = 1\),
\(x = 4\) and the
\(x\)-axis. Identify
the formula for volume of the solid of revolution obtained by revolving this region about
the \(x\)-axis.}
{\(V =\pi \int _{ 1}^{4}x\, \mathrm{d}x\)}{\(V =\int _{ 1}^{4}x\, \mathrm{d}x\)}{\(V =\pi \int _{ 1}^{4}\sqrt{x}\, \mathrm{d}x\)}{\(V =\int _{ 1}^{4}\sqrt{x}\, \mathrm{d}x\)}{}{}}
\LANGpl{
\Question{
Wykres przedstawia funkcję  \(f\colon y = \sqrt{x}\). 
Wskaż wzór na objętość bryły obrotowej, która  powstanie z obrotu obszaru ograniczonego przez wykres  funkcji  \(f\)  na przedziale 
\([ 1;\, 4] \), prostymi
\(x = 1\),
\(x = 4\), a osią 
\(x\) wokół osi \(x\).}
{\(V =\pi \int _{ 1}^{4}x\, \mathrm{d}x\)}{\(V =\int _{ 1}^{4}x\, \mathrm{d}x\)}{\(V =\pi \int _{ 1}^{4}\sqrt{x}\, \mathrm{d}x\)}{\(V =\int _{ 1}^{4}\sqrt{x}\, \mathrm{d}x\)}{}{}}
\LANGcs{
\Question{
Na obrázku je graf funkce \(f\colon y = \sqrt{x}\).
Určete vztah, podle kterého vypočítáme objem tělesa,
které vznikne rotací rovinného obrazce ohraničeného osou
\(x\), grafem
funkce \(f\) na
intervalu \(\langle 1;\, 4\rangle \) a
přímkami \(x = 1\),
\(x = 4\) kolem
osy \(x\).}
{\(V =\pi \int _{ 1}^{4}x\, \mathrm{d}x\)}{\(V =\int _{ 1}^{4}x\, \mathrm{d}x\)}{\(V =\pi \int _{ 1}^{4}\sqrt{x}\, \mathrm{d}x\)}{\(V =\int _{ 1}^{4}\sqrt{x}\, \mathrm{d}x\)}{}{}}
\LANGsk{
\Question{
Na obrázku je graf funkcie \(f\colon y = \sqrt{x}\).
Určte vzťah, podľa ktorého vypočítame objem telesa,
ktoré vznikne rotáciou rovinného obrazca ohraničeného osou
\(x\), grafom
funkcie \(f\) na
intervale \(\langle 1;\, 4\rangle \) a
priamkami \(x = 1\),
\(x = 4\) okolo
osy \(x\).}
{\(V =\pi \int _{ 1}^{4}x\, \mathrm{d}x\)}{\(V =\int _{ 1}^{4}x\, \mathrm{d}x\)}{\(V =\pi \int _{ 1}^{4}\sqrt{x}\, \mathrm{d}x\)}{\(V =\int _{ 1}^{4}\sqrt{x}\, \mathrm{d}x\)}{}{}}
\LANGes{
\Question{
En la imagen se puede ver la gráfica de la función \(f(x)= \sqrt{x}\). Consideremos la región limitada por la gráfica de la función \(f\) en el intervalo 
\([ 1;\, 4] \), las rectas \(x = 1\),
\(x = 4\) y el eje \(x\). Identifica
la fórmula del volumen del sólido de revolución obtenido por rotación de esta región alrededor del eje \(x\).}
{\(V =\pi \int _{ 1}^{4}x\, \mathrm{d}x\)}{\(V =\int _{ 1}^{4}x\, \mathrm{d}x\)}{\(V =\pi \int _{ 1}^{4}\sqrt{x}\, \mathrm{d}x\)}{\(V =\int _{ 1}^{4}\sqrt{x}\, \mathrm{d}x\)}{}{}}
\End

\Data\ID{9000100007}
\Updated{2021-08-19 10:08:33}
\Level{B}
\SubArea{Applications of definite integral}
\IMG{001000_07-figure0.pdf}
\LANGen{
\Question{
The function \(f(x)= \sqrt{x}\) 
is graphed in the picture. Consider the region bounded by the graph of
\(f\) on
\([ 1;\, 4] \), lines
\(x = 1\),
\(x = 4\) and the
\(x\)-axis.
Find the volume of the solid of revolution obtained by revolving this region about the
\(x\)-axis.}
{\(\frac{15}
 {2} \pi \)}{\(\frac{17}
 {2} \pi \)}{\(\frac{17}
 {2} \pi ^{2}\)}{\(\frac{15}
 {2} \pi ^{2}\)}{}{}}
\LANGpl{
\Question{
Wykres przedstawia funkcję \(f\colon y = \sqrt{x}\) 
Wskaż objętość bryły obrotowej, która  powstanie z obrotu obszaru ograniczonego przez wykres  funkcji \(f\)na przedziale 
\([ 1;\, 4] \), prostymi
\(x = 1\),
\(x = 4\), a osią 
\(x\) wokół osi \(x\)}
{\(\frac{15}
 {2} \pi \)}{\(\frac{17}
 {2} \pi \)}{\(\frac{17}
 {2} \pi ^{2}\)}{\(\frac{15}
 {2} \pi ^{2}\)}{}{}}
\LANGcs{
\Question{
Na obrázku je graf funkce \(f\colon y = \sqrt{x}\).
Vypočítejte objem tělesa, které vznikne rotací rovinného obrazce ohraničeného
osou \(x\), grafem
funkce \(f\) na
intervalu \(\langle 1;\, 4\rangle \) a
přímkami \(x = 1\),
\(x = 4\) kolem
osy \(x\).}
{\(\frac{15}
 {2} \pi \)}{\(\frac{17}
 {2} \pi \)}{\(\frac{17}
 {2} \pi ^{2}\)}{\(\frac{15}
 {2} \pi ^{2}\)}{}{}}
\LANGsk{
\Question{
Na obrázku je graf funkcie \(f\colon y = \sqrt{x}\).
Vypočítajte objem telesa, ktoré vznikne rotáciou rovinného obrazca ohraničeného
osou \(x\), grafom
funkcie \(f\) na
intervale \(\langle 1;\, 4\rangle \) a
priamkami \(x = 1\),
\(x = 4\) okolo
osy \(x\).}
{\(\frac{15}
 {2} \pi \)}{\(\frac{17}
 {2} \pi \)}{\(\frac{17}
 {2} \pi ^{2}\)}{\(\frac{15}
 {2} \pi ^{2}\)}{}{}}
\LANGes{
\Question{
En la imagen se puede ver la gráfica de la función \(f(x)= \sqrt{x}\). Consideremos la región limitada por la función \(f\) en
el intervalo \([ 1;\, 4] \), las rectas \(x = 1\),
\(x = 4\) y el eje \(x\).
Halla el volumen del sólido de revolución obtenido por la rotación de esta región alrededor del eje \(x\).}
{\(\frac{15}
 {2} \pi \)}{\(\frac{17}
 {2} \pi \)}{\(\frac{17}
 {2} \pi ^{2}\)}{\(\frac{15}
 {2} \pi ^{2}\)}{}{}}
\End

\Data\ID{9000100008}
\Updated{2020-08-03 04:08:28}
\Level{B}
\SubArea{Applications of definite integral}
\IMG{001000_08-figure0.pdf}
\LANGen{
\Question{
Part of the graph of the function \(f(x) = \frac{1} 
{x}\) 
is shown in the picture. Complete the following sentence: „Formula 
\[ 
 V =\pi \int _{ 1}^{2}x^{-2}\, \mathrm{d}x
\]
determines the volume of the solid of revolution obtained by revolving region
bounded by ...”}
{\(x\)-axis,
graph of \(f\) 
on \([ 1;\, 2] \) 
and lines \(x = 1\),
\(x = 2\) 
about \(x\)-axis.}{\(y\)-axis,
graph of \(f\) 
on \([ 1;\, 2] \) 
and lines \(y = 1\),
\(y = \frac{1} 
{2}\) 
about \(x\)-axis.}{\(x\)-axis,
graph of \(f^{2}\) 
on \([ 1;\, 2] \) 
and lines \(x = 1\),
\(x = 2\) 
about \(x\)-axis.}{\(y\)-axis,
graph of \(f^{2}\) 
on \([ 1;\, 2] \) 
and lines \(y = 1\),
\(y = \frac{1} 
{2}\) 
about \(x\)-axis.}{}{}}
\LANGpl{
\Question{
Wykres przedstawia część wykresu funkcji \(f\colon y = \frac{1} 
{x}\). Dokończ zdanie: „Wzór określa 
\[ 
 V =\pi \int _{ 1}^{2}x^{-2}\, \mathrm{d}x
\]
 objętość bryły obrotowej, która powstanie z obrotu obszaru ograniczonego przez}
{oś \(x\),
wykres funkcji  \(f\) 
na przedziale  \([ 1;\, 2] \) 
oraz prostymi  \(x = 1\),
\(x = 2\) 
wokół osi  \(x\).}{oś \(y\),
wykres funkcji \(f\) 
na przedziale \([ 1;\, 2] \) 
oraz prostymi \(y = 1\),
\(y = \frac{1} 
{2}\) 
wokół osi  \(x\).}{oś \(x\),
wykres funkcji\(f^{2}\) 
na przedziale \([ 1;\, 2] \) 
oraz prostymi \(x = 1\),
\(x = 2\) 
wokół osi  \(x\).}{oś \(y\),
wykres funkcji \(f^{2}\) 
na przedziale \([ 1;\, 2] \) 
oraz prostymi \(y = 1\),
\(y = \frac{1} 
{2}\) 
wokół osi \(x\).}{}{}}
\LANGcs{
\Question{
Na obrázku je část grafu funkce
\(f\colon y = \frac{1} 
{x}\).
Doplňte následující větu tak,
aby vznikl pravdivý výrok: „Objem
\(V =\pi \int _{ 1}^{2}x^{-2}\, \mathrm{d}x\) 
má těleso, které vznikne rotací
rovinného obrazce ohraničeného
...”}
{osou \(x\),
grafem funkce \(f\) 
na intervalu \(\langle 1;\, 2\rangle \) 
a přímkami \(x = 1\),
\(x = 2\) 
kolem osy \(x\).}{osou \(y\),
grafem funkce \(f\) 
na intervalu \(\langle 1;\, 2\rangle \) 
a přímkami \(y = 1\),
\(y = \frac{1} 
{2}\) 
kolem osy \(x\).}{osou \(x\),
grafem funkce \(f^{2}\) 
na intervalu \(\langle 1;\, 2\rangle \) 
a přímkami \(x = 1\),
\(x = 2\) 
kolem osy \(x\).}{osou \(y\),
grafem funkce \(f^{2}\) 
na intervalu \(\langle 1;\, 2\rangle \) 
a přímkami \(y = 1\),
\(y = \frac{1} 
{2}\) 
kolem osy \(x\).}{}{}}
\LANGsk{
\Question{
Na obrázku je časť grafu funkcie
\(f\colon y = \frac{1} 
{x}\).
Doplňte nasledujúcu vetu tak,
aby vznikol pravdivý výrok: „Objem
\(V =\pi \int _{ 1}^{2}x^{-2}\, \mathrm{d}x\) 
má teleso, ktoré vznikne rotáciou
rovinného obrazce ohraničeného
...”}
{osou \(x\),
grafom funkcie \(f\) 
na intervale \(\langle 1;\, 2\rangle \) 
a priamkami \(x = 1\),
\(x = 2\) 
okolo osy \(x\).}{osou \(y\),
grafom funkcie \(f\) 
na intervale \(\langle 1;\, 2\rangle \) 
a priamkami \(y = 1\),
\(y = \frac{1} 
{2}\) 
okolo osy \(x\).}{osou \(x\),
grafom funkcie \(f^{2}\) 
na intervale \(\langle 1;\, 2\rangle \) 
a priamkami \(x = 1\),
\(x = 2\) 
okolo osy \(x\).}{osou \(y\),
grafom funkcie \(f^{2}\) 
na intervale \(\langle 1;\, 2\rangle \) 
a priamkami \(y = 1\),
\(y = \frac{1} 
{2}\) 
okolo osy \(x\).}{}{}}
\LANGes{
\Question{
En la imagen se puede ver una parte de la gráfica de la función \(f(x) = \frac{1} 
{x}\). Completa la siguiente frase para que sea verdadera: "La fórmula 
\[ 
 V =\pi \int _{ 1}^{2}x^{-2}\, \mathrm{d}x
\]
determina el volumen del sólido de revolución obtenido por rotación de la región limitada por ...”}
{el eje \(x\),
la gráfica de la función \(f\) 
en el intervalo \([ 1;\, 2] \) 
y las rectas \(x = 1\),
\(x = 2\) 
alrededor del eje \(x\).}{el eje \(y\),
la gráfica de la función \(f\) 
en el intervalo \([ 1;\, 2] \) 
y las rectas \(y = 1\),
\(y = \frac{1} 
{2}\) 
alrededor del eje \(x\).}{el eje \(x\),
la gráfica de la función \(f^{2}\) 
en el intervalo \([ 1;\, 2] \) 
y las rectas \(x = 1\),
\(x = 2\) 
alrededor del eje \(x\).}{el eje \(y\),
la gráfica de la función \(f^{2}\) 
en el intervalo \([ 1;\, 2] \) 
y las rectas \(y = 1\),
\(y = \frac{1} 
{2}\) 
alrededor del eje \(x\).}{}{}}
\End

\Data\ID{9000100009}
\Updated{2021-08-19 10:08:58}
\Level{B}
\SubArea{Applications of definite integral}
\IMG{001000_09-figure0.pdf}
\LANGen{
\Question{
Part of the graph of the function \(f(x) = \frac{1} 
{x}\) 
is shown in the picture. Consider the region bounded by
\(x\)-axis,
graph of \(f\) 
and
lines \(x = 1\) and 
\(x = 4\). Find
the volume of the solid of revolution obtained by revolving this region about
\(x\)-axis.}
{\(\frac{3}
{4}\pi \)}{\(\frac{5}
{4}\pi \)}{\(\frac{5}
{3}\pi \)}{\(\frac{4}
{3}\pi \)}{}{}}
\LANGpl{
\Question{
Wykres przedstawia część wykresu funkcji \(f\colon y = \frac{1} 
{x}\). 
Wskaż objętość bryły obrotowej, która  powstanie z obrotu obszaru ograniczonego przez oś \(x\),  wykres  funkcji \(f\) oraz 
prostymi \(x = 1\),
\(x = 4\) wokół osi
\(x\).}
{\(\frac{3}
{4}\pi \)}{\(\frac{5}
{4}\pi \)}{\(\frac{5}
{3}\pi \)}{\(\frac{4}
{3}\pi \)}{}{}}
\LANGcs{
\Question{
Na obrázku je část grafu funkce \(f\colon y = \frac{1} 
{x}\).
Určete objem tělesa, které vznikne rotací rovinného obrazce ohraničeného osou
\(x\), grafem
funkce \(f\) a
přímkami \(x = 1\) a 
\(x = 4\) kolem
osy \(x\).}
{\(\frac{3}
{4}\pi \)}{\(\frac{5}
{4}\pi \)}{\(\frac{5}
{3}\pi \)}{\(\frac{4}
{3}\pi \)}{}{}}
\LANGsk{
\Question{
Na obrázku je časť grafu funkcie \(f\colon y = \frac{1} 
{x}\).
Určte objem telesa, ktoré vznikne rotáciou rovinného obrazca ohraničeného osou
\(x\), grafom
funkcie \(f\) a
priamkami \(x = 1\) a 
\(x = 4\) okolo
osy \(x\).}
{\(\frac{3}
{4}\pi \)}{\(\frac{5}
{4}\pi \)}{\(\frac{5}
{3}\pi \)}{\(\frac{4}
{3}\pi \)}{}{}}
\LANGes{
\Question{
En la imagen se puede ver una parte de la gráfica de la función \(f(x) = \frac{1} 
{x}\). Consideremos la región limitada por el eje \(x\),
la gráfica de \(f\) 
en el intervalo \([ 1;\, 4] \) y
las rectas \(x = 1\),
\(x = 4\). Halla el volumen del sólido de revolución obtenido por rotación de esta región alrededor del eje \(x\).}
{\(\frac{3}
{4}\pi \)}{\(\frac{5}
{4}\pi \)}{\(\frac{5}
{3}\pi \)}{\(\frac{4}
{3}\pi \)}{}{}}
\End

\Data\ID{1003124407}
\Updated{2020-07-15 03:07:12}
\Level{C}
\SubArea{Applications of definite integral}
\LANGen{
\Question{
Consider a conic section defined by $x^2-4y^2=1$. Find the volume of the solid obtained by rotating the given conic about the $x$-axis on the interval $[1;6]$.}
{$\frac{50}3\pi$}{$\frac{35}4\pi$}{$\frac{49}3\pi$}{$\frac{50}3$}{}{}}
\LANGpl{
\Question{
Dana jest krzywa stożkowa określona przez $x^2-4y^2=1$. Wskaz objętość bryły uzyskanej przez obrót danej krzywej wokół osi  $x$ w przedziale $\langle1;6\rangle$.}
{$\frac{50}3\pi$}{$\frac{35}4\pi$}{$\frac{49}3\pi$}{$\frac{50}3$}{}{}}
\LANGcs{
\Question{
Kuželosečka $x^2-4y^2=1$ vymezí rotací kolem osy $x$ na intervalu $\langle1;6\rangle$ rotační těleso. Určete jeho objem.}
{$\frac{50}3\pi$}{$\frac{35}4\pi$}{$\frac{49}3\pi$}{$\frac{50}3$}{}{}}
\LANGsk{
\Question{
Kužeľosečka $x^2-4y^2=1$ vymedzí rotáciu okolo osi $x$ na intervale $\langle1;6\rangle$ rotačné teleso. Určte jeho objem.}
{$\frac{50}3\pi$}{$\frac{35}4\pi$}{$\frac{49}3\pi$}{$\frac{50}3$}{}{}}
\LANGes{
\Question{
Considera una sección cónica definida por $x^2-4y^2=1$. Encuentra el volumen del sólido obtenido por rotación del cono alrededor del eje $x$ en el intervalo $[1;6]$.}
{$\frac{50}3\pi$}{$\frac{35}4\pi$}{$\frac{49}3\pi$}{$\frac{50}3$}{}{}}
\End

\Data\ID{1003124408}
\Updated{2020-07-15 03:07:12}
\Level{C}
\SubArea{Applications of definite integral}
\LANGen{
\Question{
Consider a conic section defined by $y^2=x-5$. Find the value of the real parameter $a$, if a solid of the volume $\frac{9\pi}2$ is obtained by rotating the given conic about the $x$-axis on the interval $[5;a]$.}
{$a=8$}{$a=2$}{$a=7.5$}{$a=9$}{}{}}
\LANGpl{
\Question{
Dana jest krzywa stożkowa określona przez $y^2=x-5$. Wyznacz wartość parametru rzeczywistego $a$, jeśli objętość bryły wynosi $\frac{9\pi}2$ , bryłę uzyskano przez obrót krzywej wokół osi $x$ w przedziale $\langle5;a\rangle$.}
{$a=8$}{$a=2$}{$a=7{,}5$}{$a=9$}{}{}}
\LANGcs{
\Question{
Určete hodnotu neznámého reálného parametru  $a$, má-li kuželosečka $y^2=x-5$ vymezit rotací kolem osy $x$ na intervalu $\langle5;a\rangle$ rotační těleso o objemu $\frac{9\pi}2$.}
{$a=8$}{$a=2$}{$a=7{,}5$}{$a=9$}{}{}}
\LANGsk{
\Question{
Určte hodnotu neznámeho reálneho parametra  $a$, ak má kužeľosečka $y^2=x-5$ vymedziť rotáciu okolo osi $x$ na intervale $\langle5;a\rangle$ rotačné teleso s objemom $\frac{9\pi}2$.}
{$a=8$}{$a=2$}{$a=7{,}5$}{$a=9$}{}{}}
\LANGes{
\Question{
Considera una sección cónica definida por $y^2=x-5$. Encuentra el valor del parámetro real $a$, si el volumen del sólido $\frac{9\pi}2$ es obtenido por rotación del cono alrededor del eje $x$ en el intervalo $[5;a]$.}
{$a=8$}{$a=2$}{$a=7.5$}{$a=9$}{}{}}
\End

\Data\ID{1003124409}
\Updated{2020-07-15 03:07:12}
\Level{C}
\SubArea{Applications of definite integral}
\LANGen{
\Question{
Consider a conic section defined by $y^2=x+a$. Find the value of the real parameter $a$, if a solid of the volume $\frac{125\pi}2$ is obtained by rotating the given conic about the $x$-axis on the interval $[0;5]$.}
{$a=10$}{$a=5$}{$a=15$}{$a=25$}{}{}}
\LANGpl{
\Question{
Dana jest krzywa stożkowa określona przez  $y^2=x+a$. Wyznacz wartość parametru rzeczywistego $a$, objętość bryły wynosi $\frac{125\pi}2$, bryłę uzyskano przez obrót krzywej wokół osi $x$ w przedziale $\langle0;5\rangle$.}
{$a=10$}{$a=5$}{$a=15$}{$a=25$}{}{}}
\LANGcs{
\Question{
Určete hodnotu neznámého reálného parametru  $a$, má-li kuželosečka  $y^2=x+a$ vymezit rotací kolem osy $x$ na intervalu $\langle0;5\rangle$ rotační těleso o objemu $\frac{125\pi}2$.}
{$a=10$}{$a=5$}{$a=15$}{$a=25$}{}{}}
\LANGsk{
\Question{
Určte hodnotu neznámeho reálneho parametra  $a$, ak má kužeľosečka  $y^2=x+a$ vymedziť rotáciu kolem osi $x$ na intervale $\langle0;5\rangle$ rotačné teleso s objemom $\frac{125\pi}2$.}
{$a=10$}{$a=5$}{$a=15$}{$a=25$}{}{}}
\LANGes{
\Question{
Considera una sección cónica definida por $y^2=x+a$. Encuentra el valor del parámetro real $a$, si el volumen del sólido $\frac{125\pi}2$ es obtenido por rotación del cono alrededor del eje $x$ en el intervalo $[0;5]$.}
{$a=10$}{$a=5$}{$a=15$}{$a=25$}{}{}}
\End

\Data\ID{1103124402}
\Updated{2020-07-22 06:07:57}
\Level{C}
\SubArea{Applications of definite integral}
\IMG{1103124402.pdf}
\LANGen{
\Question{
Find such a real number $a$, that the area of the region highlighted in green equals $6$ (see the picture).}
{$a=2$}{$a=1$}{$a=3$}{$a=4$}{}{}}
\LANGpl{
\Question{
Wskaż taką liczbę rzeczywistą $a$, aby pole powierzchni obszaru zaznaczonego kolorem zielonym było równe  $6$ (spójrz na rysunek).}
{$a=2$}{$a=1$}{$a=3$}{$a=4$}{}{}}
\LANGcs{
\Question{
Pro které reálné číslo $a$ je obsah zeleně vyznačené plochy roven $6$?}
{$a=2$}{$a=1$}{$a=3$}{$a=4$}{}{}}
\LANGsk{
\Question{
Pre ktoré reálne číslo $a$ je obsah zeleno vyznačenej plochy rovný $6$?}
{$a=2$}{$a=1$}{$a=3$}{$a=4$}{}{}}
\LANGes{
\Question{
Encuentra un número real $a$ para que el área de la parte marcada de verde sea igual a $6$ (mira la imagen).}
{$a=2$}{$a=1$}{$a=3$}{$a=4$}{}{}}
\End

\Data\ID{1103124404}
\Updated{2020-07-22 06:07:06}
\Level{C}
\SubArea{Applications of definite integral}
\IMG{1103124404.pdf}
\LANGen{
\Question{
Let a solid be obtained by rotating the blue triangle about the $x$-axis (see the picture). Find such a value of $a$, that the volume of this solid is $48\pi$.}
{$a=4$}{$a=2$}{$a=3$}{$a=6$}{}{}}
\LANGpl{
\Question{
Dana jest bryła uzyskana przez obrót niebieskiego trójkąta wokół osi $x$(spójrz na rysunek). Wyznacz taką wartość $a$, aby objętość bryły była równa $48\pi$.}
{$a=4$}{$a=2$}{$a=3$}{$a=6$}{}{}}
\LANGcs{
\Question{
Určete hodnotu reálného čísla $a$ tak, aby těleso vytvořené rotací vyznačeného modrého trojúhelníku kolem osy $x$ mělo objem $48\pi$.}
{$a=4$}{$a=2$}{$a=3$}{$a=6$}{}{}}
\LANGsk{
\Question{
Určte hodnotu reálneho čísla $a$ tak, aby teleso vytvorené rotáciou vyznačeného modrého trojuholníka okolo osi $x$ malo objem $48\pi$.}
{$a=4$}{$a=2$}{$a=3$}{$a=6$}{}{}}
\LANGes{
\Question{
Sea un sólido obtenido por rotación del triángulo azul alrededor del eje $x$ (mira la imagen). Encuentra el valor de $a$, si el volumen de este sólido es $48\pi$.}
{$a=4$}{$a=2$}{$a=3$}{$a=6$}{}{}}
\End

\Data\ID{1103124405}
\Updated{2020-07-22 06:07:02}
\Level{C}
\SubArea{Applications of definite integral}
\IMG{1103124405.pdf}
\LANGen{
\Question{
Let a solid be obtained by rotating the blue triangle about the $y$-axis (see the picture). Find the volume of this solid.}
{$32\pi$}{$32$}{$96\pi$}{$100$}{}{}}
\LANGpl{
\Question{
Dana jest bryła, która powstała przez obrót niebieskiego trójkąta wokół osi $y$ (spójrz na rysunek). Wskaż objętość bryły.}
{$32\pi$}{$32$}{$96\pi$}{$100$}{}{}}
\LANGcs{
\Question{
Jaký objem bude mít těleso vzniklé rotací vyznačeného modrého trojúhelníku kolem osy $y$?}
{$32\pi$}{$32$}{$96\pi$}{$100$}{}{}}
\LANGsk{
\Question{
Aký objem bude mať teleso vzniknuté rotáciou vyznačeného modrého trojuholníka okolo osi $y$?}
{$32\pi$}{$32$}{$96\pi$}{$100$}{}{}}
\LANGes{
\Question{
Sea un sólido obtenido por rotación del triángulo azul alrededor del eje $y$ (mira la imagen). Encuentra el volumen de este sólido.}
{$32\pi$}{$32$}{$96\pi$}{$100$}{}{}}
\End

\Data\ID{9000072901}
\Updated{2021-08-19 10:08:47}
\Level{C}
\SubArea{Applications of definite integral}
\LANGen{
\Question{
The velocity of a moving body in meters per second is given by the function
\(v(t) = 3\sqrt{t} + 2t\), where
\(t\) is a
time measured in seconds. Find the distance traveled by the body in the time interval
from \(t = 1\, \mathrm{s}\) 
to \(t = 9\, \mathrm{s}\).}
{\(132\, \mathrm{m}\)}{\(4\left (4 + \sqrt{2}\right )\mathrm{m}\)}{\(10\, \mathrm{m}\)}{}{}{}}
\LANGpl{
\Question{
Prędkość poruszającego się ciała w metrach na sekundę jest określona funkcją 
\(v(t) = 3\sqrt{t} + 2t\), gdzie 
\(t\) to czas mierzony w sekundach. Oblicz drogę poruszającego się ciała w przedziale czasu
od  \(t = 1\, \mathrm{s}\) 
do \(t = 9\, \mathrm{s}\).}
{\(132\, \mathrm{m}\)}{\(4\left (4 + \sqrt{2}\right )\mathrm{m}\)}{\(10\, \mathrm{m}\)}{}{}{}}
\LANGcs{
\Question{
Velikost okamžité rychlosti tělesa v metrech za sekundu je popsaná funkcí
\(v(t) = 3\sqrt{t} + 2t\), kde
\(t\) 
je čas v sekundách. Určete, jakou dráhu urazí těleso v době od
\(1\). do
\(9\).
sekundy.}
{\(132\, \mathrm{m}\)}{\(4\left (4 + \sqrt{2}\right )\mathrm{m}\)}{\(10\, \mathrm{m}\)}{}{}{}}
\LANGsk{
\Question{
Veľkosť okamžitej rýchlosti telesa v metroch za sekundu je popísaná funkciou
\(v(t) = 3\sqrt{t} + 2t\), kde
\(t\) 
je čas v sekundách. Určte, akú dráhu prejde teleso v dobe od
\(1\). do
\(9\).
sekundy.}
{\(132\, \mathrm{m}\)}{\(4\left (4 + \sqrt{2}\right )\mathrm{m}\)}{\(10\, \mathrm{m}\)}{}{}{}}
\LANGes{
\Question{
La velocidad de un cuerpo en movimiento viene dada por la función \(v(t) = 3\sqrt{t} + 2t\), donde el tiempo \(t\) se mide en segundos. Encuentra la distancia recorrida por el cuerpo en el intervalo temporal desde \(t = 1\, \mathrm{s}\) 
hasta  \(t = 9\, \mathrm{s}\).}
{\(132\, \mathrm{m}\)}{\(4\left (4 + \sqrt{2}\right )\mathrm{m}\)}{\(10\, \mathrm{m}\)}{}{}{}}
\End

\Data\ID{9000072902}
\Updated{2021-01-04 03:01:25}
\Level{C}
\SubArea{Applications of definite integral}
\LANGen{
\Question{
The instantaneous velocity of a moving body is proportional to the square of the time. The velocity at the
time \(t = 2\, \mathrm{s}\) 
is \(v = 6\, \mathrm{m\, s}^{-1}\).
What is the distance traveled by the body in the first \(4\) seconds?}
{\(32\, \mathrm{m}\)}{\(48\, \mathrm{m}\)}{\(24\, \mathrm{m}\)}{}{}{}}
\LANGpl{
\Question{
Prędkość poruszającego się ciała jest proporcjonalna do kwadratu czasu. Prędkość o  czasie 
 \(t = 2\, \mathrm{s}\) 
jest równa \(v = 6\, \mathrm{m\, s}^{-1}\).
Oblicz drogę poruszającego się ciała w przedziale czasu od 
\(t = 0\, \mathrm{s}\) do
\(t = 4\, \mathrm{s}\)?}
{\(32\, \mathrm{m}\)}{\(48\, \mathrm{m}\)}{\(24\, \mathrm{m}\)}{}{}{}}
\LANGcs{
\Question{
Velikost okamžité rychlosti tělesa je přímo úměrná druhé mocnině času. V čase
\(2\, \mathrm{s}\) je rychlost právě
\(6\, \mathrm{m\, s}^{-1}\). Jakou dráhu urazí
těleso za první \(4\, \mathrm{s}\)?}
{\(32\, \mathrm{m}\)}{\(48\, \mathrm{m}\)}{\(24\, \mathrm{m}\)}{}{}{}}
\LANGsk{
\Question{
Veľkosť okamžitej rýchlosti telesa je priamo úmerná druhej mocnine času. V čase
\(2\, \mathrm{s}\) je rýchlosť práve
\(6\, \mathrm{m\, s}^{-1}\). Akú dráhu urobí
teleso za prvé \(4\, \mathrm{s}\)?}
{\(32\, \mathrm{m}\)}{\(48\, \mathrm{m}\)}{\(24\, \mathrm{m}\)}{}{}{}}
\LANGes{
\Question{
La velocidad del cuerpo en movimiento está en razón directa a la segunda potencia de tiempo. La velocidad en tiempo \(t = 2\, \mathrm{s}\) 
es \(v = 6\, \mathrm{m\, s}^{-1}\).
¿Cuál es la distancia hecha por el cuerpo en el intervalo de tiempo entre \(t = 0\, \mathrm{s}\) y
\(t = 4\, \mathrm{s}\)?}
{\(32\, \mathrm{m}\)}{\(48\, \mathrm{m}\)}{\(24\, \mathrm{m}\)}{}{}{}}
\End

\Data\ID{9000072903}
\Updated{2020-07-22 01:07:09}
\Level{C}
\SubArea{Applications of definite integral}
\LANGen{
\Question{
The force required to deform a spring is proportional to the
extension of the spring. The current elongation of the spring is
\(2\, \mathrm{cm}\) and the force required to
reach this elongation is \(3\, \mathrm{N}\).
Evaluate the work required to stretch the spring from the current elongation (i.e.
\(2\, \mathrm{cm}\)) by
additional \(10\, \mathrm{cm}\).}
{\(1.05\, \mathrm{J}\)}{\(0.75\, \mathrm{J}\)}{\(0.18\, \mathrm{J}\)}{}{}{}}
\LANGpl{
\Question{
Siła potrzebna do odkształcenia sprężyny jest wprost proporcjonalna do jej przedłużenia. Aktualne wydłużenie sprężyny wynosi
\(2\, \mathrm{cm}\) siła potrzebna do osiągnięcia tego wydłużenia to 
 \(3\, \mathrm{N}\).
Oblicz  jaką pracę trzeba wykonać, aby rozciągnąć sprężynę od aktualnego wydłużenia  (tj.
\(2\, \mathrm{cm}\)) o kolejne 
 \(10\, \mathrm{cm}\).}
{\(1.05\, \mathrm{J}\)}{\(0.75\, \mathrm{J}\)}{\(0.18\, \mathrm{J}\)}{}{}{}}
\LANGcs{
\Question{
Síla nezbytně nutná k prodloužení pružiny o určitou hodnotu
je přímo úměrná tomuto prodloužení. Silou o velikosti
\(3\, \mathrm{N}\) se pružina
natáhne o \(2\, \mathrm{cm}\).
Jakou práci vykoná síla při natažení pružiny o dalších
\(10\, \mathrm{cm}\)?}
{\(1{,}05\, \mathrm{J}\)}{\(0{,}75\, \mathrm{J}\)}{\(0{,}18\, \mathrm{J}\)}{}{}{}}
\LANGsk{
\Question{
Sila nezbytne nutná k predĺženiu pružiny o určitú hodnotu
je priamo úmerná tomuto predĺženiu. Silou o veľkosti
\(3\, \mathrm{N}\) sa pružina
natiahne o \(2\, \mathrm{cm}\).
Akú prácu vykoná sila pri natiahnutí pružiny o ďalších
\(10\, \mathrm{cm}\)?}
{\(1{,}05\, \mathrm{J}\)}{\(0{,}75\, \mathrm{J}\)}{\(0{,}18\, \mathrm{J}\)}{}{}{}}
\LANGes{
\Question{
La fuerza necesaria para prolongar un resorte está en razón directa a la extensión del resorte. Para prolongar el resorte \(2\, \mathrm{cm}\) más se requiere la fuerza de \(3\, \mathrm{N}\).
Evalúa el trabajo requerido para prolongar el resorte otros \(10\, \mathrm{cm}\) más.}
{\(1.05\, \mathrm{J}\)}{\(0.75\, \mathrm{J}\)}{\(0.18\, \mathrm{J}\)}{}{}{}}
\End

\Data\ID{9000072904}
\Updated{2021-01-04 03:01:22}
\Level{C}
\SubArea{Applications of definite integral}
\LANGen{
\Question{
The force of the repulsion of two charged particles is 
\[ 
 F(x) = \frac{c} 
{x^{2}},
\]
where \(x\) is the distance
in meters and \(c\) 
a positive constant. Find the work required to increase the distance between the particles
from \(3\, \mathrm{m}\) 
to \(1\, \mathrm{m}\).}
{\(\frac{2}
{3}c\, \mathrm{J}\)}{\(\frac{1}
{3}c\, \mathrm{J}\)}{\(c\, \mathrm{J}\)}{}{}{}}
\LANGpl{
\Question{
Siła odpychania dwóch naładowanych cząstek wynosi
\[ 
 F(x) = \frac{c} 
{x^{2}},
\]
gdzie  \(x\) to odległość w metrach,
a  \(c\) to wielkość stała dodatnia.
Oblicz jaką pracę należy wykonać, aby zmienić  odległość pomiędzy cząstkami z
 \(3\, \mathrm{m}\) 
do \(1\, \mathrm{m}\).}
{\(\frac{2}
{3}c\, \mathrm{J}\)}{\(\frac{1}
{3}c\, \mathrm{J}\)}{\(c\, \mathrm{J}\)}{}{}{}}
\LANGcs{
\Question{
Dvě nabité částice se odpuzují silou, jejíž velikost v newtonech je
popsaná funkcí 
\[ 
 F(x) = \frac{c} 
{x^{2}},
\]
kde \(x\) je vzdálenost
částic v metrech a \(c\) 
nějaká kladná konstanta. Jakou práci vykonáme při přemístění částic ze
vzdálenosti \(3\, \mathrm{m}\) do
vzdálenosti \(1\, \mathrm{m}\) 
od sebe?}
{\(\frac{2}
{3}c\, \mathrm{J}\)}{\(\frac{1}
{3}c\, \mathrm{J}\)}{\(c\, \mathrm{J}\)}{}{}{}}
\LANGsk{
\Question{
Dve nabité častice sa odpudzujú silou, ktorých veľkosť v newtonoch je
popísaná funkciou 
\[ 
 F(x) = \frac{c} 
{x^{2}},
\]
kde \(x\) je vzdialenosť
častíc v metroch a \(c\) 
nejaká kladná konštanta. Akú prácu vykonáme pri premiestnení častíc zo
vzdialenosti \(3\, \mathrm{m}\) do
vzdialenosti \(1\, \mathrm{m}\) 
od seba?}
{\(\frac{2}
{3}c\, \mathrm{J}\)}{\(\frac{1}
{3}c\, \mathrm{J}\)}{\(c\, \mathrm{J}\)}{}{}{}}
\LANGes{
\Question{
Dos partículas igualmente cargadas se repelen con una fuerza en newton. La fuerza se define por la siguiente función 
\[ 
 F(x) = \frac{c} 
{x^{2}},
\]
donde \(x\) es la distancia en metros y \(c\) una constante positiva. Encuentra el trabajo requerido para desplazar las partículas de la distancia de \(3\, \mathrm{m}\) 
a la distancia de \(1\, \mathrm{m}\) entre ellas.}
{\(\frac{2}
{3}c\, \mathrm{J}\)}{\(\frac{1}
{3}c\, \mathrm{J}\)}{\(c\, \mathrm{J}\)}{}{}{}}
\End

\Data\ID{9000072905}
\Updated{2020-07-30 05:07:30}
\Level{C}
\SubArea{Applications of definite integral}
\LANGen{
\Question{
The \(1000\, \mathrm{kg}\) heavy satellite is
transported to the orbit \(150\, \mathrm{km}\) 
above the ground. Find the mechanical work required for this transport. The mass of the Earth
is \(M = 6\cdot 10^{24}\, \mathrm{kg}\), gravitational
constant \(\kappa = 6.67\cdot 10^{-11}\, \mathrm{N\, m}^{2}\mathrm{kg}^{-2}\) and Earth
radius \(R = 6\: 370\, \mathrm{km}\). Round your
result to nearest \(\mathrm{MJ}\).}
{\(1\: 445\, \mathrm{MJ}\)}{\(1\: 471\, \mathrm{MJ}\)}{\(1\: 412\, \mathrm{MJ}\)}{}{}{}}
\LANGpl{
\Question{
\(1000\, \mathrm{kg}\) satelita została wystrzelona 
na orbitę  \(150\, \mathrm{km}\) 
nad Ziemią. Oblicz pracę mechaniczną potrzebną do wystrzelenia satelity. Masa Ziemi wynosi  \(M = 6\cdot 10^{24}\, \mathrm{kg}\), stała grawitacji 
\(\kappa = 6.67\cdot 10^{-11}\, \mathrm{N\, m}^{2}\mathrm{kg}^{-2}\), promień Ziemi jest równy
 \(R = 6\: 370\, \mathrm{km}\). Zaokrągli wynik do pełnych
\(\mathrm{MJ}\).}
{\(1\: 445\, \mathrm{MJ}\)}{\(1\: 471\, \mathrm{MJ}\)}{\(1\: 412\, \mathrm{MJ}\)}{}{}{}}
\LANGcs{
\Question{
Určete práci vykonanou při vynesení tunové družice do výšky
\(150\, \mathrm{km}\) ze zemského povrchu.
Hmotnost Země je \(M = 6\cdot 10^{24}\, \mathrm{kg}\),
gravitační konstanta \(\kappa = 6{,}67\cdot 10^{-11}\, \mathrm{N\, m}^{2}\mathrm{kg}^{-2}\) 
a poloměr Země \(R = 6\: 370\, \mathrm{km}\).
Zaokrouhlete na \(\mathrm{MJ}\).}
{\(1\: 445\, \mathrm{MJ}\)}{\(1\: 471\, \mathrm{MJ}\)}{\(1\: 412\, \mathrm{MJ}\)}{}{}{}}
\LANGsk{
\Question{
Určte prácu vykonanú pri vynesení tonovej družice do výšky
\(150\, \mathrm{km}\) zo zemského povrchu.
Hmotnosť Zem je \(M = 6\cdot 10^{24}\, \mathrm{kg}\),
gravitačná konštanta \(\kappa = 6{,}67\cdot 10^{-11}\, \mathrm{N\, m}^{2}\mathrm{kg}^{-2}\) 
a polomer Zeme \(R = 6\: 370\, \mathrm{km}\).
Zaokrúhlite na \(\mathrm{MJ}\).}
{\(1\: 445\, \mathrm{MJ}\)}{\(1\: 471\, \mathrm{MJ}\)}{\(1\: 412\, \mathrm{MJ}\)}{}{}{}}
\LANGes{
\Question{
El satélite que pesa \(1000\, \mathrm{kg}\) se transportó a la órbita \(150\, \mathrm{km}\) 
sobre la Tierra. Busca el trabajo mecánico requerido para este transporte. La masa de la Tierra es \(M = 6\cdot 10^{24}\, \mathrm{kg}\), constante gravitacional
\(\kappa = 6.67\cdot 10^{-11}\, \mathrm{N\, m}^{2}\mathrm{kg}^{-2}\) y el radio de la Tierra \(R = 6\: 370\, \mathrm{km}\). Aproxima el resultado a \(\mathrm{MJ}\).}
{\(1\: 445\, \mathrm{MJ}\)}{\(1\: 471\, \mathrm{MJ}\)}{\(1\: 412\, \mathrm{MJ}\)}{}{}{}}
\End

\Data\ID{9000072906}
\Updated{2020-07-30 05:07:08}
\Level{C}
\SubArea{Applications of definite integral}
\LANGen{
\Question{
The reservoir in the form of a box is filled with the water. The vertical side of the box is
\(50\, \mathrm{cm}\) height
and \(40\, \mathrm{cm}\) 
long. Find the total force which acts on this side. The mass density of the water is
\(1\: 000\, \mathrm{kg\, m}^{-3}\) and the standard
acceleration of gravity is \(g = 9.81\, \mathrm{m\, s}^{-2}\).}
{\(490.5\, \mathrm{N}\)}{\(981\, \mathrm{N}\)}{\(245.25\, \mathrm{N}\)}{}{}{}}
\LANGpl{
\Question{
Zbiornik w kształcie pudełka jest wypełniony wodą. Pionowa strona zbiornika  ma 
\(50\, \mathrm{cm}\) wysokości
i \(40\, \mathrm{cm}\) 
szerokości. Oblicz całkowitą siłę działającą na ten bok. Gęstość wody wynosi
\(1\: 000\, \mathrm{kg\, m}^{-3}\), standardowe przyspieszenie grawitacyjne
to \(g = 9.81\, \mathrm{m\, s}^{-2}\).}
{\(490.5\, \mathrm{N}\)}{\(981\, \mathrm{N}\)}{\(245.25\, \mathrm{N}\)}{}{}{}}
\LANGcs{
\Question{
Akvárium tvaru kvádru je zcela zaplněno vodou. Jakou
silou působí voda na jeho boční stěnu, která je vysoká
\(50\, \mathrm{cm}\) a dlouhá
\(40\, \mathrm{cm}\)? Hustota vody
je \(1\: 000\, \mathrm{kg\, m}^{-3}\) a tíhové
zrychlení \(g = 9{,}81\, \mathrm{m\, s}^{-2}\).}
{\(490{,}5\, \mathrm{N}\)}{\(981\, \mathrm{N}\)}{\(245{,}25\, \mathrm{N}\)}{}{}{}}
\LANGsk{
\Question{
Akvárium tvaru kvádra je celkom zaplnené vodou. Akou
silou pôsobí voda na jeho bočnú stenu, ktorá je vysoká
\(50\, \mathrm{cm}\) a dlhá
\(40\, \mathrm{cm}\)? Hustota vody
je \(1\: 000\, \mathrm{kg\, m}^{-3}\) a tiažové
zrýchlenie \(g = 9{,}81\, \mathrm{m\, s}^{-2}\).}
{\(490{,}5\, \mathrm{N}\)}{\(981\, \mathrm{N}\)}{\(245{,}25\, \mathrm{N}\)}{}{}{}}
\LANGes{
\Question{
El acuario de forma rectangular está completamente lleno de agua. El lado vertical mide \(50\, \mathrm{cm}\) de altura y
\(40\, \mathrm{cm}\) de longitud. Halla la fuerza total que actúa sobre este lado. La densidad de agua es \(1\: 000\, \mathrm{kg\, m}^{-3}\) y la gravedad de la Tierra es \(g = 9.81\, \mathrm{m\, s}^{-2}\).}
{\(490.5\, \mathrm{N}\)}{\(981\, \mathrm{N}\)}{\(245.25\, \mathrm{N}\)}{}{}{}}
\End

\Data\ID{9000072907}
\Updated{2020-08-04 09:08:32}
\Level{C}
\SubArea{Applications of definite integral}
\LANGen{
\Question{
A homogeneous cube with the side \(10\, \mathrm{cm}\) 
is in the water. The bottom side is parallel to the water surface
\(10\, \mathrm{cm}\) below
the surface. Find the work required to move the cube to the position when the
bottom side just touches the surface of the water. The mass density of the cube is
\(2\: 000\, \mathrm{kg\, m}^{-3}\), the mass density of the
water is \(1\: 000\, \mathrm{kg\, m}^{-3}\) and the standard
acceleration of gravity is \(g = 10\, \mathrm{m\, s}^{-2}\).}
{\(1.5\, \mathrm{J}\)}{\(2\, \mathrm{J}\)}{\(1\, \mathrm{J}\)}{}{}{}}
\LANGpl{
\Question{
Sześcian jednorodny o boku \(10\, \mathrm{cm}\) 
jest znużony w wodzie. Jego dolny bok jest równoległy do powierzchni wody i znajduje się 
\(10\, \mathrm{cm}\) poniżej powierzchni wody. Oblicz jaką pracę należy wykonać, aby przesunąć  sześcian tak, aby jego dolny bok stykał się z powierzchnią wody. Gęstość sześcianu wynosi
\(2\: 000\, \mathrm{kg\, m}^{-3}\), gęstość wody 
jest równa \(1\: 000\, \mathrm{kg\, m}^{-3}\), a standardowe przyspieszenie grawitacyjne 
to  \(g = 10\, \mathrm{m\, s}^{-2}\).}
{\(1.5\, \mathrm{J}\)}{\(2\, \mathrm{J}\)}{\(1\, \mathrm{J}\)}{}{}{}}
\LANGcs{
\Question{
Homogenní krychle o hraně \(10\, \mathrm{cm}\) 
a hustotě \(2\: 000\, \mathrm{kg\, m}^{-3}\) 
je ponořena do vody tak, že její dolní stěna je
rovnoběžná s volnou hladinou vody a nachází se v hloubce
\(10\, \mathrm{cm}\).
Vypočítejte práci potřebnou ke zvednutí krychle do polohy, v níž bude
její spodní stěna právě na volné hladině kapaliny. Hustota vody je
\(1\: 000\, \mathrm{kg\, m}^{-3}\) a tíhové
zrychlení \(g = 10\, \mathrm{m\, s}^{-2}\).}
{\(1{,}5\, \mathrm{J}\)}{\(2\, \mathrm{J}\)}{\(1\, \mathrm{J}\)}{}{}{}}
\LANGsk{
\Question{
Homogénna kocka s hranou \(10\, \mathrm{cm}\) 
a hustotou \(2\: 000\, \mathrm{kg\, m}^{-3}\) 
je ponorená do vody tak, že jej dolná stena je
rovnobežná s voľnou hladinou vody a nachádza sa v hĺbke
\(10\, \mathrm{cm}\).
Vypočítajte prácu potrebnú k zdvihnutiu kocky do polohy, v ktorej bude
jej spodná stena práve na voľnej hladine kvapaliny. Hustota vody je
\(1\: 000\, \mathrm{kg\, m}^{-3}\) a tiažové
zrýchlenie \(g = 10\, \mathrm{m\, s}^{-2}\).}
{\(1{,}5\, \mathrm{J}\)}{\(2\, \mathrm{J}\)}{\(1\, \mathrm{J}\)}{}{}{}}
\LANGes{
\Question{
Un cubo homogéneo de \(10\, \mathrm{cm}\) de lado que tiene una densidad de   
\(2\: 000\, \mathrm{kg\, m}^{-3}\) está en el agua. El lado inferior es paralelo con la superficie del agua y está \(10\, \mathrm{cm}\) bajo la superficie. Halla el trabajo requerido para mover el cubo a la posición donde el lado inferior flote justo en la superficie. La densidad del agua es \(1\: 000\, \mathrm{kg\, m}^{-3}\) y la gravedad de la Tierra es \(g = 10\, \mathrm{m\, s}^{-2}\).}
{\(1.5\, \mathrm{J}\)}{\(2\, \mathrm{J}\)}{\(1\, \mathrm{J}\)}{}{}{}}
\End

\Data\ID{9000072908}
\Updated{2020-08-04 09:08:54}
\Level{C}
\SubArea{Applications of definite integral}
\LANGen{
\Question{
A \(100\, \mathrm{kg}\) heavy anchor is attached
to a \(20\, \mathrm{m}\) long rope. One meter
of the rope weights \(1\, \mathrm{kg}\).
Find the work required to raise the anchor
\(20\, \mathrm{m}\) higher. The standard
acceleration of gravity is \(9.81\, \mathrm{m\, s}^{-2}\).
Neglect the buoyancy (the force from the Archimedes law).}
{\(21\: 582\, \mathrm{J}\)}{\(23\: 544\, \mathrm{J}\)}{\(19\: 620\, \mathrm{J}\)}{}{}{}}
\LANGpl{
\Question{
\(100\, \mathrm{kg}\) kotwica jest doczepiona do 
 \(20\, \mathrm{m}\) liny. Metr liny waży
\(1\, \mathrm{kg}\).
Oblicz jaką pracę należy wykonać, aby podnieść kotwicę 
\(20\, \mathrm{m}\) do góry. Standardowe przyspieszenie grawitacyjne wynosi
 \(9.81\, \mathrm{m\, s}^{-2}\).
Nie uwzględniamy siły wyporu.}
{\(21\: 582\, \mathrm{J}\)}{\(23\: 544\, \mathrm{J}\)}{\(19\: 620\, \mathrm{J}\)}{}{}{}}
\LANGcs{
\Question{
Kotva o hmotnosti \(100\, \mathrm{kg}\) visí na
kotevním laně délky \(20\, \mathrm{m}\). Jeden
metr lana má hmotnost \(1\, \mathrm{kg}\).
Jakou práci vykonáme, jestliže zvedneme kotvu o
\(20\, \mathrm{m}\)? Tíhové
zrychlení je \(9{,}81\, \mathrm{m\, s}^{-2}\).
Vztlakovou sílu zanedbáme.}
{\(21\: 582\, \mathrm{J}\)}{\(23\: 544\, \mathrm{J}\)}{\(19\: 620\, \mathrm{J}\)}{}{}{}}
\LANGsk{
\Question{
Kotva s hmotnosťou \(100\, \mathrm{kg}\) visí na
kotevnom lane dĺžky \(20\, \mathrm{m}\). Jeden
meter lana má hmotnosť \(1\, \mathrm{kg}\).
Akú prácu vykonáme, ak zdvihneme kotvu o
\(20\, \mathrm{m}\)? Tiažové
zrýchlenie je \(9{,}81\, \mathrm{m\, s}^{-2}\).
Vztlakovú silu neberieme do úvahy.}
{\(21\: 582\, \mathrm{J}\)}{\(23\: 544\, \mathrm{J}\)}{\(19\: 620\, \mathrm{J}\)}{}{}{}}
\LANGes{
\Question{
Un ancla de peso \(100\, \mathrm{kg}\) está colgada en una cuerda que mide \(20\, \mathrm{m}\). Un metro de esta cuerda pesa \(1\, \mathrm{kg}\).
Busca el trabajo requerido para subir el ancla \(20\, \mathrm{m}\). La gravedad de la Tierra es \(9.81\, \mathrm{m\, s}^{-2}\).
Desatendamos la flotación (la fuerza de principio de Arquímedes).}
{\(21\: 582\, \mathrm{J}\)}{\(23\: 544\, \mathrm{J}\)}{\(19\: 620\, \mathrm{J}\)}{}{}{}}
\End
